\documentclass[UTF8,12pt,a4paper,fontset=fandol]{ctexbook}

% 支持从项目根目录或本笔记目录编译。
\makeatletter
\def\input@path{{../}}
\makeatother
\usepackage{researchnotes}

\title{抽象代数笔记}
\author{}
\date{}

\begin{document}

\frontmatter
\maketitle
\tableofcontents

\chapter{绪论与全书大纲}

\section{研究对象与课程定位}

整数可以相加、相乘，几何变换可以复合，矩阵可以相乘，多项式可以作带余除法。这些运算的具体对象不同，却反复呈现结合律、单位元、逆元和分配律等共同规律。\keyterm{抽象代数（abstract algebra）}从这些规律出发，研究带有运算的集合，以及保持运算的映射。

本书围绕三类问题组织内容：一个代数结构由哪些元素生成；哪些映射能够保持它的运算；怎样借助子结构、商结构和分解定理理解整体。\keyterm{群（group）}用于研究单一可逆运算与对称性，\keyterm{环（ring）}和\keyterm{域（field）}用于研究加法、乘法与方程，\keyterm{模（module）}则把线性空间中的标量推广到环，连接群论、环论和线性代数。

课程按本科两学期的完整学习范围设计，预备知识为集合与映射、整数整除和基础线性代数。必要的证明方法与整数知识在正文中补足。群、环、域构成核心主线；模论、标准形及标为选学的专题用于深化。微积分不是主要先修课程，研究生层次的交换代数、表示论和同调代数不在本书系统展开的范围内。

本书按上述路线逐章展开。第\ref{ch:preliminaries}至\ref{ch:cyclic-permutation}章已写成正文，给出定义、证明与例题；其余章节暂保留写作大纲，以说明后续内容及依赖关系。大纲中列出的证明任务不视为已经完成的论证。

\section{全书课程大纲}

\begin{description}
  \item[第一篇：预备知识与群论] 从数学语言和对称性出发，建立群的基本工具，再研究有限群的结构。
\end{description}
\begin{enumerate}
  \item 第\ref{ch:preliminaries}章：数学语言与预备知识。
  \item 第\ref{ch:groups}章：群、子群与生成。
  \item 第\ref{ch:cyclic-permutation}章：循环群与置换群。
  \item 第\ref{ch:cosets}章：陪集、指数与拉格朗日定理。
  \item 第\ref{ch:quotient-groups}章：正规子群、商群与同态定理。
  \item 第\ref{ch:actions}章：群作用与计数。
  \item 第\ref{ch:sylow}章：西罗定理与有限群。
  \item 第\ref{ch:abelian}章：直积与有限阿贝尔群。
  \item 第\ref{ch:solvable-groups}章：群列、单群与可解群。
\end{enumerate}

\begin{description}
  \item[第二篇：环与多项式] 把整数与多项式中的整除规律放入统一框架，以理想构造商环，并为扩域作准备。
\end{description}
\begin{enumerate}[resume]
  \item 第\ref{ch:rings}章：环、整环、域与环同态。
  \item 第\ref{ch:ideals}章：理想、商环与中国剩余定理。
  \item 第\ref{ch:factorization}章：整除性与唯一分解。
  \item 第\ref{ch:polynomials}章：多项式环与不可约性。
\end{enumerate}

\begin{description}
  \item[第三篇：域与伽罗瓦理论] 从添加方程的根出发，用域扩张描述代数数，再用根的对称性研究方程能否用根式求解。
\end{description}
\begin{enumerate}[resume]
  \item 第\ref{ch:extensions}章：域扩张与代数元素。
  \item 第\ref{ch:constructions}章：尺规作图与二次扩张。
  \item 第\ref{ch:splitting}章：分裂域、可分扩张与正规扩张。
  \item 第\ref{ch:finite-fields}章：有限域。
  \item 第\ref{ch:galois}章：伽罗瓦对应与伽罗瓦群计算。
  \item 第\ref{ch:radicals}章：根式扩张与方程的可解性。
\end{enumerate}

\begin{description}
  \item[第四篇：模与线性代数的统一] 把阿贝尔群和线性变换放入模的语言，利用结构定理解释标准形。
\end{description}
\begin{enumerate}[resume]
  \item 第\ref{ch:modules}章：模、子模与模同态。
  \item 第\ref{ch:pid-modules}章：主理想整环上的有限生成模与标准形。
\end{enumerate}

\section{学习路线与知识依赖}

初学时先完成第\ref{ch:preliminaries}至\ref{ch:actions}章，掌握定义、例子、同态和商结构的基本方法，再学习西罗定理与有限群分解。环论从第\ref{ch:rings}章开始，主要依赖前面的群论基础；第\ref{ch:sylow}至\ref{ch:solvable-groups}章的较深内容可以与环论交替学习。

域论依赖第\ref{ch:polynomials}章的不可约多项式与商环构造，以及线性代数中的基、维数和线性无关。第\ref{ch:constructions}章是域扩张的独立应用，略过它不影响后续理论。第\ref{ch:radicals}章同时依赖可解群与伽罗瓦对应。模论在完成环论后即可开始，不以伽罗瓦理论为先修。

一学期入门路线可以选取第\ref{ch:preliminaries}至\ref{ch:quotient-groups}章、群作用的基本部分、第\ref{ch:rings}至\ref{ch:extensions}章，并以有限域或尺规作图作结。两学期路线覆盖完整核心主线，再按学习进度加入模论与选学内容。具体进度应以能独立完成定义验证、计算与基础证明为准。

\section{统一记号与学习要求}

整数、有理数、实数、复数分别记为 $\mathbb Z,\mathbb Q,\mathbb R,\mathbb C$；自然数约定为 $\mathbb N=\{0,1,2,\ldots\}$。群通常记为 $G$，单位元记为 $e$，子群关系记为 $H\leq G$；加法群的单位元记为 $0$。整数 $n\geq 1$ 的剩余类集合记为 $\mathbb Z/n\mathbb Z$，整数 $a$ 的剩余类记为 $\overline a$。映射复合 $f\circ g$ 表示先作用 $g$，再作用 $f$。

环均指有乘法单位元的结合环，允许零环；环同态与子环均要求保持单位元。整环和域额外要求 $1\ne0$。环不默认交换，进入交换环的结论时会明确说明。模默认指含幺左模。\keyterm{同态（homomorphism）}保持相应运算，\keyterm{同构（isomorphism）}则是可逆的结构保持映射；同构记为 $\cong$，不与集合相等混用。

每章展开时采用“问题与例子、定义、基本性质、关键定理及证明、应用、习题”的顺序。习题包含定义辨析、计算、证明和综合应用；反例用于说明假设不可随意删除。正文中正式引入符号与术语，各章大纲中的“练习与自检”说明后续题组应检验的能力。

\mainmatter
\part{预备知识与群论}

\chapter{数学语言与预备知识}\label{ch:preliminaries}

抽象代数把具体计算中的共同规律写成公理，再由公理推导结论。学习这一过程，需要准确区分对象与表示、存在与唯一、命题与逆命题。本章准备三项工具：用映射描述结构，用等价关系构造新对象，用整数整除提供最初的模型。所需知识限于中学代数与集合的基本概念。

\section{命题、量词与证明方法}\label{sec:logic}

\keyterm{命题（proposition）}是具有确定真值的陈述。对于命题 $P,Q$，$P\Rightarrow Q$ 表示 $P$ 是 $Q$ 的充分条件、$Q$ 是 $P$ 的必要条件；$P\Leftrightarrow Q$ 表示两者互为充要条件。证明等价命题通常需要分别证明两个方向。逆命题 $Q\Rightarrow P$ 不由原命题推出，而逆否命题 $\neg Q\Rightarrow\neg P$ 与原命题等价。

\keyterm{量词（quantifier）}指定陈述适用的范围。设 $P(x)$ 是关于 $x\in X$ 的条件，则 $\forall x\in X,\ P(x)$ 要求每个元素满足条件，$\exists x\in X,\ P(x)$ 只要求至少一个元素满足条件。否定时应同时交换量词并否定条件：
\[
  \neg(\forall x\in X,\ P(x))\ \Longleftrightarrow\
  \exists x\in X,\ \neg P(x),\qquad
  \neg(\exists x\in X,\ P(x))\ \Longleftrightarrow\
  \forall x\in X,\ \neg P(x).
\]
记号 $\exists!x\in X,\ P(x)$ 表示存在唯一的满足条件的元素。证明它时，先构造一个满足条件的元素，再证明任意两个满足条件的元素相等。

\begin{example}[量词顺序]
在实数中，$\forall x\in\mathbb R,\ \exists y\in\mathbb R,\ x+y=0$ 成立：对给定的 $x$ 取 $y=-x$。但 $\exists y\in\mathbb R,\ \forall x\in\mathbb R,\ x+y=0$ 不成立，因为分别取 $x=0,1$ 会要求同一个 $y$ 同时等于 $0$ 和 $-1$。后文“每个元素有逆元”中的逆元可以依赖于该元素。
\end{example}

直接证明从假设推出结论；反证法在假设结论不成立后导出矛盾；反例则给出一个满足假设但不满足结论的对象。对有限或可逐步构造的对象，常用归纳法。它的基础可以表述为自然数的\keyterm{良序原理（well-ordering principle）}：每个非空的 $\mathbb N$ 的子集都有最小元。本书将此作为自然数的基本性质。

\begin{proposition}[数学归纳法]\label{prop:induction}
设 $P(n)$ 是关于 $n\in\mathbb N$ 的命题。若 $P(0)$ 成立，并且对每个 $n\in\mathbb N$ 都有 $P(n)\Rightarrow P(n+1)$，则所有 $P(n)$ 均成立。
\end{proposition}
\begin{proof}
若结论不成立，反例指标组成非空集合，取其最小元 $m$。由 $P(0)$ 成立知 $m>0$；最小性给出 $P(m-1)$，归纳假设继而给出 $P(m)$，矛盾。
\end{proof}

同一论证还给出强归纳法：证明 $P(n)$ 时，可以假设所有较小指标对应的命题都成立。使用归纳法必须写明起始指标与归纳推进；“已经验证很多情形”不能替代对任意指标的论证。

\section{集合、映射与运算}\label{sec:maps-operations}

集合由其元素确定，证明 $A=B$ 可以分别证明 $A\subseteq B$ 和 $B\subseteq A$。集合 $X,Y$ 的\keyterm{笛卡尔积（Cartesian product）}为 $X\times Y=\{(x,y):x\in X,\ y\in Y\}$，其中有序对相等当且仅当对应分量相等。有限集合 $X$ 的元素个数记为 $|X|$。

\begin{definition}[映射与像]
\keyterm{映射（map）}$f:X\to Y$ 给每个 $x\in X$ 指定唯一元素 $f(x)\in Y$。$X$ 称为定义域，$Y$ 称为陪域。对 $A\subseteq X$、$B\subseteq Y$，定义\keyterm{像（image）}与\keyterm{原像（preimage）}
\[
  f(A)=\{f(x):x\in A\},\qquad
  f^{-1}(B)=\{x\in X:f(x)\in B\}.
\]
整个映射的像记为 $\operatorname{im}f=f(X)$，它可以小于陪域 $Y$。
\end{definition}

若 $f(x)=f(x')$ 总能推出 $x=x'$，则 $f$ 是\keyterm{单射（injection）}；若 $f(X)=Y$，则 $f$ 是\keyterm{满射（surjection）}；同时为单射、满射的映射称为\keyterm{双射（bijection）}。原像 $f^{-1}(B)$ 对任意映射都有意义，不要求 $f$ 存在逆映射。

给定 $f:X\to Y$、$g:Y\to Z$，其复合定义为 $(g\circ f)(x)=g(f(x))$。恒等映射 $\operatorname{id}_X:X\to X$ 满足 $\operatorname{id}_X(x)=x$。复合满足结合律，因为两种括号方式在每个元素上的取值都是 $h(g(f(x)))$，但一般不满足交换律，甚至交换次序后的复合可能没有定义。

\begin{proposition}[双射与逆映射]\label{prop:inverse-map}
映射 $f:X\to Y$ 为双射，当且仅当存在映射 $g:Y\to X$ 满足 $g\circ f=\operatorname{id}_X$、$f\circ g=\operatorname{id}_Y$；这样的 $g$ 唯一，记为 $f^{-1}$。
\end{proposition}
\begin{proof}
若 $f$ 为双射，则对每个 $y\in Y$ 都有唯一的 $x\in X$ 满足 $f(x)=y$，令 $g(y)=x$ 即得所需映射。反之，若两式成立，由 $f(x)=f(x')$ 两边作用 $g$ 得 $x=x'$，且每个 $y=f(g(y))$ 都属于像，故 $f$ 为双射。若 $g,h$ 都满足两式，则 $g=g\circ(f\circ h)=(g\circ f)\circ h=h$。
\end{proof}

\begin{definition}[二元运算]
非空集合 $X$ 上的\keyterm{二元运算（binary operation）}是映射 $*:X\times X\to X$，通常将 $*(a,b)$ 写成 $a*b$。若对任意 $a,b,c\in X$ 有 $(a*b)*c=a*(b*c)$，称其满足结合律；若对任意 $a,b\in X$ 有 $a*b=b*a$，称其满足交换律。
\end{definition}

运算的结果仍在 $X$ 中这一要求称为封闭性。例如减法是 $\mathbb Z$ 上的二元运算，却不是 $\mathbb N$ 上的二元运算；它在 $\mathbb Z$ 上不满足结合律，因为 $(1-1)-1=-1$ 而 $1-(1-1)=1$。同一集合配上不同运算会得到不同结构，因此谈论代数对象时必须同时说明集合和运算。

\section{等价关系、划分与商集}\label{sec:equivalence-quotients}

在许多问题中，只关心某些差别而忽略另一些差别。例如整数除以 $n$ 时，可以只记录余数。这种“视为相同”的规则必须保持一致，才能把原来的元素合并成新对象。

\begin{definition}[等价关系与商集]
集合 $X$ 上的关系 $\sim$ 若满足自反性 $x\sim x$、对称性 $x\sim y\Rightarrow y\sim x$、传递性 $x\sim y,\ y\sim z\Rightarrow x\sim z$，则称为\keyterm{等价关系（equivalence relation）}。元素 $x$ 的\keyterm{等价类（equivalence class）}为 $[x]=\{y\in X:y\sim x\}$，全部等价类组成\keyterm{商集（quotient set）}$X/{\sim}=\{[x]:x\in X\}$。
\end{definition}

\begin{proposition}[等价关系与划分]\label{prop:equivalence-partition}
等价关系的等价类都是非空的，覆盖 $X$，且任意两个等价类或者相等，或者不相交。反过来，将 $X$ 分成若干两两不交的非空子集，就唯一确定一个以这些子集为等价类的等价关系。
\end{proposition}
\begin{proof}
由 $x\sim x$，每个 $x$ 都在 $[x]$ 中。若 $z\in[x]\cap[y]$，则 $z\sim x$、$z\sim y$，从而 $x\sim y$。对任意 $u\in[x]$，由 $u\sim x\sim y$ 得 $u\in[y]$，故 $[x]\subseteq[y]$；交换 $x,y$ 得到相等。反过来，定义 $x\sim y$ 当且仅当两者属于同一个分块，逐项满足三条性质，且分块已经唯一规定了哪些元素相关。
\end{proof}

\begin{example}[整数同余]\label{ex:integer-congruence}
固定正整数 $n$，定义 $a\equiv b\pmod n$ 当且仅当 $a-b=nk$ 对某个 $k\in\mathbb Z$ 成立。由 $a-a=0$、$b-a=-(a-b)$ 及 $a-c=(a-b)+(b-c)$，这是等价关系。等价类记为 $\overline a$，商集记为 $\mathbb Z/n\mathbb Z$。例如模 $3$ 时，$\overline1$ 包含 $\ldots,-5,-2,1,4,7,\ldots$；$1$ 只是该类的一个代表元，不是这个类本身。
\end{example}

\begin{proposition}[通过商集定义映射]\label{prop:quotient-set-map}
设 $\sim$ 是 $X$ 上的等价关系，$f:X\to Y$。存在映射 $\overline f:X/{\sim}\to Y$ 满足 $\overline f([x])=f(x)$，当且仅当 $x\sim x'$ 蕴含 $f(x)=f(x')$；存在时此映射唯一。
\end{proposition}
\begin{proof}
若 $\overline f$ 存在，同一个等价类只能有一个像，故条件必要。若条件成立，则 $f(x)$ 不依赖 $[x]$ 的代表元选择，所给公式确定一个映射。每个等价类都形如 $[x]$，因而公式也唯一确定了所有取值。
\end{proof}

上述代表元无关性称为\keyterm{良定义性（well-definedness）}。类似地，若要由 $X$ 上的运算 $*$ 定义 $[x]\mathbin{\overline *}[y]=[x*y]$，必须且只需证明
\begin{equation}\label{eq:operation-descends}
  x\sim x',\quad y\sim y'\quad\Longrightarrow\quad x*y\sim x'*y'.
\end{equation}

\begin{example}[剩余类运算的良定义性]\label{ex:residue-operations}
在 $\mathbb Z/n\mathbb Z$ 上令 $\overline a+\overline b=\overline{a+b}$、$\overline a\,\overline b=\overline{ab}$。若 $a'=a+nu$、$b'=b+nv$，其中 $u,v\in\mathbb Z$，则
\[
  (a'+b')-(a+b)=n(u+v),\qquad
  a'b'-ab=n(av+bu+nuv).
\]
两种结果都不依赖代表元，因此加法、乘法均良定义。反之，在 $\mathbb Z$ 上定义 $a\sim b$ 当且仅当 $|a|=|b|$，虽然这是等价关系，公式 $[a]+[b]=[a+b]$ 却不良定义：$[1]=[-1]$，但 $[1+1]=[2]\ne[0]=[-1+1]$。等价关系本身并不保证运算能传到商集。
\end{example}

\section{整数整除与贝祖等式}\label{sec:integer-divisibility}

对整数 $a,b$，若存在 $c\in\mathbb Z$ 使 $b=ac$，就说 $a$\keyterm{整除（divide）}$b$，记为 $a\mid b$。若 $a\mid b$ 且 $a\mid c$，则 $a$ 整除 $ub+vc$，其中 $u,v$ 为任意整数。这一简单性质把整除问题与整数线性组合联系起来。

\begin{theorem}[带余除法]\label{thm:integer-division}
对任意 $a\in\mathbb Z$ 与正整数 $b$，存在唯一整数 $q,r$ 满足 $a=bq+r$ 且 $0\leq r<b$。
\end{theorem}
\begin{proof}
集合 $S=\{a-bq:q\in\mathbb Z,\ a-bq\geq0\}$ 非空，例如可取 $q=-|a|$。取其最小元 $r=a-bq$。若 $r\geq b$，则 $r-b=a-b(q+1)$ 仍在 $S$ 中，却小于 $r$，矛盾，故 $0\leq r<b$。若另有 $a=bq'+r'$ 满足同样条件，则 $b(q-q')=r'-r$，右端绝对值小于 $b$，只能为 $0$，于是 $q=q'$、$r=r'$。
\end{proof}

由此每个模 $n$ 剩余类都恰有一个代表元属于 $\{0,1,\ldots,n-1\}$，所以 $|\mathbb Z/n\mathbb Z|=n$。这里的唯一性是指选定这个代表元范围之后的唯一性。

\begin{definition}[最大公因数与互素]
对不全为零的整数 $a,b$，它们的\keyterm{最大公因数（greatest common divisor）}$\gcd(a,b)$ 是最大的正公因数。若 $\gcd(a,b)=1$，称 $a,b$\keyterm{互素（coprime）}。
\end{definition}

\begin{theorem}[贝祖等式]\label{thm:integer-bezout}
设 $a,b\in\mathbb Z$ 不全为零，$d=\gcd(a,b)$。存在 $u,v\in\mathbb Z$ 使
\begin{equation}\label{eq:integer-bezout}
  d=ua+vb.
\end{equation}
并且每个公因数都整除 $d$。
\end{theorem}
\begin{proof}
所有正整数线性组合 $ua+vb$ 构成非空集合，取其最小元 $d_0$。用带余除法写 $a=qd_0+r$，其中 $0\leq r<d_0$。由于 $d_0$ 是 $a,b$ 的整数线性组合，$r=a-qd_0$ 也是；若 $r>0$ 就违背最小性，故 $r=0$。同理 $d_0\mid b$。任一公因数 $c$ 都整除 $d_0$，所以每个正公因数不超过 $d_0$。于是 $d_0=d$，并得到所需表示及整除性质。
\end{proof}

\begin{corollary}[互素消去与欧几里得引理]\label{cor:coprime-cancellation}
若 $\gcd(a,b)=1$ 且 $a\mid bc$，则 $a\mid c$。特别地，若素数 $p$ 整除 $bc$，则 $p\mid b$ 或 $p\mid c$；这里素数指大于 $1$、正因数只有 $1$ 与自身的整数。
\end{corollary}
\begin{proof}
由 $ua+vb=1$ 得 $c=uac+vbc$，右端两项都被 $a$ 整除。若 $p\nmid b$，素数的定义给出 $\gcd(p,b)=1$，将前一结论用于 $p,b,c$ 即得 $p\mid c$。
\end{proof}

\begin{theorem}[整数的唯一素因子分解]\label{thm:integer-factorization}
每个整数 $n>1$ 都能写成有限个素数的乘积；除因子的排列顺序外，这种分解唯一。
\end{theorem}
\begin{proof}
对 $n$ 用强归纳法。若 $n$ 是素数，结论成立；否则 $n=ab$，其中 $1<a,b<n$，分别分解 $a,b$ 即得存在性。若 $p_1\cdots p_r=q_1\cdots q_s$ 是两种素数分解，推论\ref{cor:coprime-cancellation}反复用于右端可知 $p_1$ 整除某个 $q_j$，从而 $p_1=q_j$。消去这一对相等因子并继续；不可能一边先耗尽而另一边仍有素数因子，因为后者的乘积大于 $1$。因此因子及其重数相同。
\end{proof}

\keyterm{欧几里得算法（Euclidean algorithm）}计算最大公因数：先取 $|a|,|b|$；若第二项为 $0$，输出第一项，否则反复用“除数、余数”替换当前两项。其正确性来自 $\gcd(a,b)=\gcd(b,a-qb)$，因为两对整数有相同的公因数；非零余数严格递减，保证有限步终止。最后一个非零余数就是最大公因数，反向代入各步等式可得到式\eqref{eq:integer-bezout}的系数。

\begin{example}[求最大公因数与贝祖系数]\label{ex:euclidean-algorithm}
计算 $\gcd(252,198)$ 并求一组贝祖系数。
\end{example}
\begin{solve}
连续除法给出
\[
  252=198+54,\quad 198=3\cdot54+36,\quad
  54=36+18,\quad 36=2\cdot18.
\]
故最大公因数为 $18$。反向代入得到
\[
  18=54-36=4\cdot54-198=4\cdot252-5\cdot198.
\]
系数为 $u=4,v=-5$；贝祖系数一般不唯一，但最大公因数唯一。
\end{solve}

\begin{theorem}[线性同余方程]\label{thm:linear-congruence}
设 $n\geq1$，$a,b\in\mathbb Z$，$d=\gcd(a,n)$。方程 $ax\equiv b\pmod n$ 有整数解当且仅当 $d\mid b$。若有解 $x_0$，则全部解为
\begin{equation}\label{eq:linear-congruence-solutions}
  x=x_0+k\frac nd,\qquad k\in\mathbb Z,
\end{equation}
在模 $n$ 意义下恰有 $d$ 个不同的解。
\end{theorem}
\begin{proof}
若 $ax-b=ny$，则 $d\mid b$。反之，由 $ua+vn=d$ 且 $d\mid b$，取 $x_0=u(b/d)$ 即有 $ax_0\equiv b\pmod n$。任意两个解的差满足 $n\mid a(x-x_0)$；除去 $d$ 后，由 $\gcd(a/d,n/d)=1$ 和推论\ref{cor:coprime-cancellation}知 $n/d\mid x-x_0$，得到式\eqref{eq:linear-congruence-solutions}，反向代入也确为解。该式中两个参数给出相同剩余类当且仅当它们之差被 $d$ 整除，因此可取 $k=0,\ldots,d-1$。
\end{proof}

\section{例题、练习与自检}

\begin{example}[求解线性同余]
求 $14x\equiv8\pmod{30}$ 的全部整数解。
\end{example}
\begin{solve}
有 $\gcd(14,30)=2\mid8$，方程可约为 $7x\equiv4\pmod{15}$。由 $7\cdot13\equiv1\pmod{15}$ 得 $x\equiv52\equiv7\pmod{15}$。全部整数解为 $x=7+15k$，$k\in\mathbb Z$；模 $30$ 的两个解是 $\overline7,\overline{22}$。
\end{solve}

\begin{enumerate}
  \item\label{exer:quantifier-uniqueness} 写出“对每个 $x\in X$，存在唯一 $y\in Y$ 使 $P(x,y)$ 成立”的量词表达式及其否定。说明证明存在性与唯一性各需要完成什么任务。
  \item\label{exer:composition-injective-surjective} 设 $f:X\to Y$、$g:Y\to Z$。证明：若 $g\circ f$ 为单射，则 $f$ 为单射；若 $g\circ f$ 为满射，则 $g$ 为满射。分别给出另一个映射不具有相应性质的例子。
  \item\label{exer:image-preimage} 对 $f:X\to Y$ 与 $A\subseteq X$，证明 $A\subseteq f^{-1}(f(A))$；证明对每个 $A$ 都取等号当且仅当 $f$ 单射。
  \item\label{exer:quotient-square-map} 在 $\mathbb Z$ 上定义 $a\sim b$ 当且仅当 $a-b$ 为偶数。求商集，并证明 $\overline a\mapsto\overline{a^2}$ 是该商集到自身的良定义映射。
  \item\label{exer:absolute-value-quotient} 验证“具有相同绝对值”是 $\mathbb Z$ 上的等价关系。例\ref{ex:residue-operations}说明加法不能传到商集；判断乘法能否传到商集，并证明你的结论。
  \item\label{exer:bezout-computation} 用欧几里得算法计算 $\gcd(391,299)$，再求一组贝祖系数。
  \item\label{exer:linear-congruences} 求解 $18x\equiv12\pmod{30}$；判断 $18x\equiv7\pmod{30}$ 是否有解。
  \item\label{exer:gcd-lcm-product} 设 $a,b$ 为正整数，证明 $\gcd(a,b)\operatorname{lcm}(a,b)=ab$，其中最小公倍数 $\operatorname{lcm}(a,b)$ 指最小的正公倍数。可使用定理\ref{thm:integer-factorization}。
\end{enumerate}

\textbf{部分提示与结果。} 第\ref{exer:composition-injective-surjective}题均可取 $X=Z=\{0\}$、$Y=\{0,1\}$，令 $f$ 为包含映射、$g$ 为常值映射。第\ref{exer:image-preimage}题的逆向证明可取单元素子集。第\ref{exer:absolute-value-quotient}题利用 $|ab|=|a||b|$。第\ref{exer:bezout-computation}题结果为 $23=-3\cdot391+4\cdot299$。第\ref{exer:linear-congruences}题第一个方程的解为 $x\equiv4\pmod5$，第二个无解。第\ref{exer:gcd-lcm-product}题比较各素因子在两边的指数。

本章的关键检查顺序是：先确定对象与量词，再验证定义是否有意义，最后证明所需性质。尤其在商集上，必须先解决代表元无关性，才有资格讨论映射或运算的进一步性质。

\chapter{群、子群与生成}\label{ch:groups}

正三角形的旋转与反射可以复合，每个变换可以撤销，恒等变换不改变图形。这些性质不依赖三角形的具体大小，也出现在可逆矩阵与双射的复合中。群的公理抽取的正是“具有单位元的可逆结合运算”。本章先从公理推出基本运算法则，再研究保留同一运算的子结构及其生成方式。

\section{从对称性与可逆运算到群}\label{sec:group-axioms}

\begin{definition}[群]
非空集合 $G$ 配上二元运算 $(a,b)\mapsto ab$，若满足下列条件，就称为群：
\begin{enumerate}
  \item 对任意 $a,b,c\in G$，有 $(ab)c=a(bc)$；
  \item 存在 $e\in G$，使任意 $a\in G$ 都满足 $ea=ae=a$；
  \item 对每个 $a\in G$，存在 $b\in G$，使 $ab=ba=e$。
\end{enumerate}
元素 $e$ 称为\keyterm{单位元（identity element）}，$b$ 称为 $a$ 的\keyterm{逆元（inverse element）}。若还满足 $ab=ba$ 对任意 $a,b\in G$ 成立，则称为\keyterm{阿贝尔群（abelian group）}，亦称交换群。
\end{definition}

群的\keyterm{阶（order）}指集合 $G$ 的基数，记为 $|G|$；有限阶群称为有限群。只含单位元的群称为平凡群。结合律使有限乘积不依赖加括号的方式，因此通常省略括号；它并不允许任意交换因子的次序。

\begin{proposition}[单位元与逆元的唯一性]\label{prop:group-uniqueness}
群的单位元唯一，每个元素的逆元唯一。
\end{proposition}
\begin{proof}
若 $e,e'$ 都是单位元，则 $e=ee'=e'$。固定单位元，若 $b,c$ 都是 $a$ 的逆元，则 $b=be=b(ac)=(ba)c=ec=c$。
\end{proof}

因此可将 $a$ 的逆元记为 $a^{-1}$。加法记号下，单位元写成 $0$，逆元写成 $-a$；这种记号通常用于阿贝尔群。

\begin{proposition}[基本运算法则]\label{prop:group-calculation}
设 $a,b,c$ 为群 $G$ 的元素，则
\[
  (a^{-1})^{-1}=a,\qquad (ab)^{-1}=b^{-1}a^{-1},\qquad e^{-1}=e.
\]
并且有左、右消去律：$ab=ac\Rightarrow b=c$，$ba=ca\Rightarrow b=c$。方程 $ax=b$ 与 $ya=b$ 分别有唯一解 $x=a^{-1}b$、$y=ba^{-1}$。
\end{proposition}
\begin{proof}
$a$ 是 $a^{-1}$ 的逆元，$e$ 是自身的逆元，由唯一性得到相应等式。又
\[
  (ab)(b^{-1}a^{-1})=a(bb^{-1})a^{-1}=e,\qquad
  (b^{-1}a^{-1})(ab)=b^{-1}(a^{-1}a)b=e,
\]
故得到乘积的逆元公式。对 $ab=ac$ 左乘 $a^{-1}$ 得 $b=c$；右消去同理。解方程时分别左乘、右乘逆元得到候选解，代回验证存在性，消去律保证唯一性。
\end{proof}

\begin{example}[加法群与乘法群]
$(\mathbb Z,+)$、$(\mathbb Q,+)$、$(\mathbb R,+)$ 都是阿贝尔群。由例\ref{ex:residue-operations}，$\mathbb Z/n\mathbb Z$ 上的加法良定义；结合律来自整数加法，单位元为 $\overline0$，$\overline a$ 的逆元为 $\overline{-a}$，所以它也是阿贝尔群。非零实数 $\mathbb R^\times=\mathbb R\setminus\{0\}$ 在乘法下为阿贝尔群，但 $\mathbb Z\setminus\{0\}$ 在乘法下不是群，因为 $2$ 没有整数乘法逆元。
\end{example}

\begin{example}[双射的复合]
集合 $X$ 的全部双射组成集合 $\operatorname{Sym}(X)$。复合在此集合内封闭且满足结合律，单位元为 $\operatorname{id}_X$，命题\ref{prop:inverse-map}给出每个双射的逆。因此 $\operatorname{Sym}(X)$ 在复合下为群。这一构造保证几何对称与置换可以用同一套群论语言研究。
\end{example}

\section{子群与子群判别法}\label{sec:subgroups}

\begin{definition}[子群]
群 $G$ 的非空子集 $H$，若在 $G$ 的运算限制下也构成群，就称为 $G$ 的\keyterm{子群（subgroup）}，记为 $H\leq G$。$\{e\}$ 与 $G$ 总是子群；若 $H\ne G$，则称为真子群。
\end{definition}

子群的单位元与逆元必与原群一致。事实上，若 $h_0$ 是子群单位元，则 $h_0h_0=h_0=h_0e$，在 $G$ 中消去左因子得 $h_0=e$；逆元的一致性再由命题\ref{prop:group-uniqueness}得到。

\begin{theorem}[子群判别法]\label{thm:subgroup-test}
设 $G$ 为群，$H\subseteq G$。则 $H\leq G$ 当且仅当 $H\ne\varnothing$，并且对任意 $a,b\in H$ 都有 $ab^{-1}\in H$。
\end{theorem}
\begin{proof}
必要性由子群对逆元和乘法封闭得到。反之，先取 $a\in H$，得 $e=aa^{-1}\in H$。对任意 $b\in H$，得 $b^{-1}=eb^{-1}\in H$；于是对 $a,b\in H$，有 $ab=a(b^{-1})^{-1}\in H$。结合律直接继承自 $G$，故 $H$ 为群。
\end{proof}

非空条件不能省略，因为空集也形式上满足“对任意 $a,b\in H$”之后的条件。加法群中，判别条件写成 $a-b\in H$。

\begin{proposition}[有限子集的判别法]\label{prop:finite-subgroup-test}
群 $G$ 的非空有限子集 $H$ 为子群，当且仅当 $H$ 对乘法封闭。
\end{proposition}
\begin{proof}
只需证明充分性。固定 $a\in H$，左乘映射 $L_a:H\to H$、$x\mapsto ax$ 由消去律为单射。有限集合上的单射必为满射，故存在 $x\in H$ 使 $ax=a$，从而 $x=e$，所以 $e\in H$。再由 $L_a$ 满射，存在 $b\in H$ 使 $ab=e$，于是 $b=a^{-1}$。对所有 $a\in H$ 都成立，结合乘法封闭即得子群。
\end{proof}

有限性不可删除：$\mathbb N\subseteq(\mathbb Z,+)$ 包含 $0$ 且对加法封闭，但不包含 $1$ 的加法逆元。另一个常用事实是子群的交仍为子群。

\begin{proposition}[子群的交]\label{prop:subgroup-intersection}
同一个群 $G$ 的任意一族子群的交仍为子群，其中空族的交约定为 $G$。
\end{proposition}
\begin{proof}
非空族中每个子群都包含 $e$，故交非空；若 $a,b$ 同时属于所有子群，则 $ab^{-1}$ 也同时属于它们。应用定理\ref{thm:subgroup-test}即可。空族的情形由约定成立。
\end{proof}

\section{生成子群与元素的阶}\label{sec:generated-subgroups}

对 $a\in G$，定义 $a^0=e$，正整数次幂为相应个 $a$ 的乘积，并对 $n>0$ 令 $a^{-n}=(a^{-1})^n$。对任意整数 $m,n$，有
\begin{equation}\label{eq:group-power-laws}
  a^ma^n=a^{m+n},\qquad (a^m)^n=a^{mn}.
\end{equation}
第一式在指数同为非负或同为非正时由合并乘积得到，异号时逐对消去 $a$ 与 $a^{-1}$ 得到；第二式在 $n>0$ 时反复用第一式，在 $n<0$ 时再取逆元。式\eqref{eq:group-power-laws}只涉及同一个元素，不允许在非交换群中直接写 $(ab)^n=a^nb^n$。

\begin{definition}[生成子群]
对 $S\subseteq G$，定义由 $S$\keyterm{生成的子群（generated subgroup）}为
\[
  \langle S\rangle=\bigcap_{S\subseteq H\leq G}H.
\]
若 $\langle S\rangle=G$，称 $S$ 为 $G$ 的生成集。单元素集生成的子群记为 $\langle a\rangle$；若一个群由一个元素生成，称为\keyterm{循环群（cyclic group）}。
\end{definition}

由于 $G$ 本身包含 $S$，上式所取的子群族非空。由命题\ref{prop:subgroup-intersection}，$\langle S\rangle$ 是子群，并且是包含 $S$ 的最小子群。这里最小是指包含于其他每个包含 $S$ 的子群。

\begin{proposition}[生成子群的元素描述]\label{prop:generated-words}
$\langle S\rangle$ 恰由所有有限乘积 $s_1^{\varepsilon_1}\cdots s_k^{\varepsilon_k}$ 组成，其中 $k\geq0$、$s_i\in S$、$\varepsilon_i\in\{1,-1\}$；$k=0$ 时约定空乘积为 $e$。特别地，$\langle a\rangle=\{a^m:m\in\mathbb Z\}$，且 $\langle\varnothing\rangle=\{e\}$。
\end{proposition}
\begin{proof}
记所有这类乘积组成的集合为 $W$。拼接乘积表明 $W$ 对乘法封闭，逆元公式表明取逆相当于逆序排列并改变各指数符号，故仍在 $W$ 中；又 $e\in W$，所以 $W$ 为包含 $S$ 的子群。反过来，任何包含 $S$ 的子群必须包含这些有限乘积，故 $W$ 等于所有这些子群的交。
\end{proof}

\begin{definition}[元素的阶]
若存在正整数 $d$ 使 $a^d=e$，就把其中最小者称为 $a$ 的阶，记为 $\operatorname{ord}(a)$；若不存在，则称 $a$ 为无限阶元素，记为 $\operatorname{ord}(a)=\infty$。
\end{definition}

\begin{proposition}[幂相等的判别]\label{prop:order-divisibility}
若 $\operatorname{ord}(a)=d<\infty$，则对整数 $k,m,n$ 有
\[
  a^k=e\ \Longleftrightarrow\ d\mid k,\qquad
  a^m=a^n\ \Longleftrightarrow\ m\equiv n\pmod d.
\]
此时 $\langle a\rangle=\{e,a,\ldots,a^{d-1}\}$，这些元素两两不同。若 $a$ 无限阶，则所有整数次幂两两不同。
\end{proposition}
\begin{proof}
若 $d\mid k$，用幂法则即得 $a^k=e$。反之，写 $k=qd+r$，其中 $0\leq r<d$，则 $e=a^k=a^r$；由 $d$ 的最小性必有 $r=0$。再将 $a^m=a^n$ 化成 $a^{m-n}=e$，得到其余有限阶结论。无限阶时，若 $a^m=a^n$ 且 $m\ne n$，则某个正整数次幂为 $e$，与定义矛盾。
\end{proof}

因此元素的阶恰是其生成子群的阶。有限群中的每个元素都有有限阶，但目前尚不能据此断言其阶整除群阶；后一结论需要第\ref{ch:cosets}章的拉格朗日定理。

\section{基本例子与反例库}\label{sec:basic-group-examples}

\begin{example}[矩阵群]
对 $n\geq1$，全体可逆 $n$ 阶实矩阵组成\keyterm{一般线性群（general linear group）}$\operatorname{GL}_n(\mathbb R)$。乘积可逆，因为 $(AB)^{-1}=B^{-1}A^{-1}$；单位元为单位矩阵 $I_n$，结合律继承自矩阵乘法。行列式等于 $1$ 的矩阵组成子群 $\operatorname{SL}_n(\mathbb R)$，称为\keyterm{特殊线性群（special linear group）}，因为 $\det(AB^{-1})=\det A/\det B=1$。
\end{example}

当 $n\geq2$ 时，$\operatorname{GL}_n(\mathbb R)$ 不是阿贝尔群。例如在 $n=2$ 时，取
\[
  A=\begin{pmatrix}1&1\\0&1\end{pmatrix},\qquad
  B=\begin{pmatrix}1&0\\1&1\end{pmatrix},
\]
则 $AB=\left(\begin{smallmatrix}2&1\\1&1\end{smallmatrix}\right)$，$BA=\left(\begin{smallmatrix}1&1\\1&2\end{smallmatrix}\right)$，二者不等。对 $n>2$，取以这两个矩阵及 $I_{n-2}$ 为对角块的分块对角矩阵，即得同样的反例。矩阵集合在加法下的群结构与这里的乘法群结构不同，不能混用单位元和逆元。

\begin{example}[二面体群]\label{ex:dihedral-group}
正 $n$ 边形（$n\geq3$）的全部刚性对称组成\keyterm{二面体群（dihedral group）}$D_{2n}$。将顶点标为 $\mathbb Z/n\mathbb Z$，旋转 $r$ 与一个反射 $s$ 在顶点上的作用为 $r(k)=k+1$、$s(k)=-k$，所有计算均模 $n$。从作用直接得到
\begin{equation}\label{eq:dihedral-relations}
  r^n=e,\qquad s^2=e,\qquad sr=r^{-1}s.
\end{equation}
全部对称为 $r^a$ 或 $r^as$，其中 $0\leq a<n$：一个对称由某顶点的像和相邻顶点的像确定，至多有 $2n$ 种，而这 $2n$ 个变换两两不同并且都是对称。因此 $|D_{2n}|=2n$，且 $D_{2n}=\langle r,s\rangle$。
\end{example}

利用式\eqref{eq:dihedral-relations}，可将任意乘积化为统一形式：
\begin{equation}\label{eq:dihedral-multiplication}
  (r^as^b)(r^cs^d)=r^{a+(-1)^bc}s^{b+d},
  \qquad b,d\in\{0,1\},
\end{equation}
其中 $r$ 的指数模 $n$，$s$ 的指数模 $2$。特别地，$(r^as)^2=e$，所以每个反射的阶为 $2$。关系式来自已经构造出的变换群，因此不会把形式符号的运算误当作未经验证的群。

\begin{example}[四元数群]\label{ex:quaternion-group}
令 $i\in\mathbb C$ 满足 $i^2=-1$，取复矩阵
\[
  P=\begin{pmatrix}i&0\\0&-i\end{pmatrix},\qquad
  Q=\begin{pmatrix}0&1\\-1&0\end{pmatrix}.
\]
直接相乘得到 $P^2=Q^2=-I_2$、$QP=-PQ$。集合
\[
  Q_8=\{\pm I_2,\pm P,\pm Q,\pm PQ\}
\]
含八个不同的可逆矩阵；上述关系把任意乘积化为 $\pm P^aQ^b$，其中 $a,b\in\{0,1\}$，故它对乘法封闭。由命题\ref{prop:finite-subgroup-test}，这是群，称为\keyterm{四元数群（quaternion group）}。常将 $P,Q,PQ$ 分别记为 $\boldsymbol i,\boldsymbol j,\boldsymbol k$；矩阵构造同时保证了结合律。
\end{example}

\section{例题、练习与自检}

\begin{example}[整数加法群的全部子群]\label{ex:integer-subgroups}
证明 $(\mathbb Z,+)$ 的每个子群都唯一地形如 $m\mathbb Z=\{mk:k\in\mathbb Z\}$，其中 $m\geq0$。
\end{example}
\begin{solve}
先由差的封闭性验证 $m\mathbb Z$ 是子群。反过来，设 $H\leq\mathbb Z$。若 $H=\{0\}$，取 $m=0$；否则 $H$ 包含某个非零数及其相反数，因而包含正数。取其中最小的正数 $m$。对子群中任意 $a$ 作带余除法 $a=qm+r$，$0\leq r<m$；由于 $r=a-qm\in H$，最小性迫使 $r=0$，故 $H\subseteq m\mathbb Z$。另一包含关系由 $m\in H$ 得到。非零子群的最小正数唯一，平凡子群对应 $m=0$，故表示唯一。
\end{solve}

\begin{example}[区分两个八阶群]
在 $D_8$ 中，$r,r^3$ 的阶为 $4$，$r^2$ 与四个反射的阶为 $2$；在 $Q_8$ 中，只有 $-I_2$ 的阶为 $2$，其余六个非单位元的阶均为 $4$。因此两个群虽有相同阶，却不能通过保持乘法的双射相互对应：这样的双射保持每个等式 $a^k=e$，从而保持元素阶及其分布。
\end{example}

\begin{enumerate}
  \item\label{exer:exponent-two-group} 设群 $G$ 的每个元素都满足 $a^2=e$。证明 $G$ 为阿贝尔群。逆命题是否成立？
  \item\label{exer:subgroup-union} 设 $H,K\leq G$，证明 $H\cup K$ 是子群当且仅当 $H\subseteq K$ 或 $K\subseteq H$。
  \item\label{exer:generated-minimality} 设 $S\subseteq G$、$H\leq G$。证明 $S\subseteq H$ 当且仅当 $\langle S\rangle\subseteq H$，并说明 $S$ 不必对子群运算封闭。
  \item\label{exer:dihedral-products} 在 $D_{10}$ 中化简 $(r^2s)(r^4s)$ 与 $sr^3s$，并求每种形式元素的逆元。
  \item\label{exer:quaternion-relations} 证明 $Q_8=\langle P,Q\rangle$，求 $PQ$ 的逆，并验证 $PQP^{-1}=Q^{-1}$。
  \item\label{exer:integer-two-generators} 求 $(\mathbb Z,+)$ 中由 $18$ 与 $30$ 生成的子群，并推广到由任意不全为零的整数 $a,b$ 生成的情形。
  \item\label{exer:conjugate-order} 设 $a\in G$ 为有限阶元素。证明 $a$ 与 $a^{-1}$ 阶相同，并证明 $gag^{-1}$ 与 $a$ 阶相同，其中 $g\in G$。
\end{enumerate}

\textbf{部分提示与结果。} 第\ref{exer:exponent-two-group}题比较 $(ab)^{-1}$ 的两种表达。第\ref{exer:subgroup-union}题若两个子群互不包含，取 $h\in H\setminus K$、$k\in K\setminus H$，分析 $hk$ 属于哪一边会导致矛盾。第\ref{exer:dihedral-products}题两式分别为 $r^3$、$r^2$；反射均为自身的逆。第\ref{exer:integer-two-generators}题答案为 $6\mathbb Z$，一般结果为 $\gcd(a,b)\mathbb Z$，使用贝祖等式。第\ref{exer:conjugate-order}题先证明 $(gag^{-1})^k=ga^kg^{-1}$。

\chapter{循环群与置换群}\label{ch:cyclic-permutation}

循环群的每个元素都是同一元素的幂，因而其结构可以转化为整数的整除问题。置换群则直接记录对象之间的重新排列，是有限非交换群的基本模型。本章分别建立循环群的完整分类与置换的标准计算方法，并通过子群、元素阶和奇偶性识别两类结构的特征。

\section{循环群及其分类}\label{sec:cyclic-classification}

由上一章，循环群形如 $G=\langle g\rangle=\{g^k:k\in\mathbb Z\}$，元素 $g$ 称为一个\keyterm{生成元（generator）}。幂法则给出 $g^mg^n=g^{m+n}=g^{n+m}=g^ng^m$，所以每个循环群都是阿贝尔群；逆命题将在本章用反例否定。

\begin{definition}[群同构]\label{def:group-isomorphism}
设 $G,H$ 为群。若双射 $f:G\to H$ 满足 $f(ab)=f(a)f(b)$ 对所有 $a,b\in G$ 成立，则称 $f$ 为群同构，并记 $G\cong H$。
\end{definition}

同构保留单位元与逆元：由 $f(e)f(e)=f(e)$ 消去得 $f(e)=e$，再由 $f(a)f(a^{-1})=e$ 得 $f(a^{-1})=f(a)^{-1}$；这里两边的 $e$ 分别属于相应的群。逆双射也保持乘法，因为若 $y=f(a),z=f(b)$，则 $f^{-1}(yz)=ab=f^{-1}(y)f^{-1}(z)$。因此同构表示运算结构相同，而不仅是元素数量相同。一般不要求双射的结构保持映射将在第\ref{ch:quotient-groups}章系统讨论。

\begin{theorem}[循环群的分类]\label{thm:cyclic-classification}
每个无限循环群都同构于 $(\mathbb Z,+)$；每个 $n$ 阶循环群都同构于 $(\mathbb Z/n\mathbb Z,+)$。
\end{theorem}
\begin{proof}
设 $G=\langle g\rangle$。若 $g$ 无限阶，映射 $k\mapsto g^k$ 从 $\mathbb Z$ 到 $G$ 满射，且由命题\ref{prop:order-divisibility}为单射；幂法则保证它把整数加法对应到群乘法。若 $|G|=n<\infty$，同一命题给出 $\operatorname{ord}(g)=n$，故 $\overline k\mapsto g^k$ 从 $\mathbb Z/n\mathbb Z$ 到 $G$ 良定义且双射，并同样保持运算。
\end{proof}

分类定理中的同构取决于生成元的选择，一般不唯一。它说明研究任意循环群时可以改用整数或剩余类计算。

\begin{theorem}[循环群的子群]\label{thm:cyclic-subgroups}
循环群的每个子群仍是循环群。更具体地，若 $H\leq\langle g\rangle$ 且 $H\ne\{e\}$，取使 $g^d\in H$ 的最小正整数 $d$，则 $H=\langle g^d\rangle$。
\end{theorem}
\begin{proof}
平凡子群由 $e$ 生成。若 $H$ 非平凡，取 $g^k\in H$ 且 $g^k\ne e$；必要时取逆元，可以找到正指数，所以所述 $d$ 存在。对任意 $g^m\in H$，写 $m=qd+r$，其中 $0\leq r<d$，则 $g^r=g^m(g^d)^{-q}\in H$。若 $r>0$，与 $d$ 最小矛盾，故 $r=0$，于是 $g^m\in\langle g^d\rangle$。另一包含关系由 $g^d\in H$ 得到。
\end{proof}

\begin{proposition}[循环群中的元素阶]\label{prop:cyclic-element-order}
若 $\operatorname{ord}(g)=n<\infty$，则对任意整数 $k$，有
\begin{equation}\label{eq:cyclic-element-order}
  \operatorname{ord}(g^k)=\frac{n}{\gcd(n,k)}.
\end{equation}
若 $g$ 无限阶，则 $k\ne0$ 时 $g^k$ 仍为无限阶。
\end{proposition}
\begin{proof}
有限阶时，$(g^k)^m=e$ 等价于 $n\mid km$。记 $c=\gcd(n,k)$，约去 $c$ 并利用 $\gcd(n/c,k/c)=1$，可知这又等价于 $n/c\mid m$。最小的正整数 $m$ 因而为 $n/c$。无限阶时，若 $(g^k)^m=e$ 且 $m>0$，则 $km=0$，与 $k\ne0$ 矛盾。
\end{proof}

\section{生成元与子群格}\label{sec:cyclic-subgroup-lattice}

\begin{corollary}[生成元判别]\label{cor:cyclic-generators}
设 $G=\langle g\rangle$ 为 $n$ 阶循环群。则 $g^k$ 生成 $G$ 当且仅当 $\gcd(n,k)=1$。无限循环群 $\langle g\rangle$ 恰有两个生成元 $g,g^{-1}$。
\end{corollary}
\begin{proof}
有限情形中，$g^k$ 生成 $G$ 当且仅当 $\langle g^k\rangle$ 的阶等于 $n$，使用式\eqref{eq:cyclic-element-order}即可。无限情形中，若 $g\in\langle g^k\rangle$，则 $g=g^{km}$ 对某个整数 $m$ 成立；由无限阶性有 $km=1$，所以 $k=\pm1$，两者也确实生成全群。
\end{proof}

\keyterm{欧拉函数（Euler's totient function）}$\varphi(n)$ 定义为 $1\leq k\leq n$ 中与 $n$ 互素的整数个数。故 $n$ 阶循环群恰有 $\varphi(n)$ 个生成元。约定 $\varphi(1)=1$，与平凡群的唯一元素生成该群相符。

\begin{theorem}[有限循环群的子群分类]\label{thm:finite-cyclic-subgroups}
设 $G=\langle g\rangle$ 的阶为 $n$。对每个正因子 $d\mid n$，$G$ 恰有一个 $d$ 阶子群
\begin{equation}\label{eq:cyclic-subgroup-divisor}
  H_d=\langle g^{n/d}\rangle.
\end{equation}
这些子群包含了 $G$ 的全部子群，并满足 $H_d\subseteq H_c$ 当且仅当 $d\mid c$。
\end{theorem}
\begin{proof}
式\eqref{eq:cyclic-element-order}给出 $|H_d|=d$。若 $H$ 非平凡，按定理\ref{thm:cyclic-subgroups}写成 $\langle g^m\rangle$，其中 $m$ 是使 $g^m\in H$ 的最小正整数。对 $n$ 除以 $m$，由 $g^n=e\in H$ 与同一定理的余数论证得到 $m\mid n$，故 $|H|=n/m$，从阶即可唯一恢复 $m$。平凡子群对应 $d=1$。

若 $d\mid c$，则 $g^{n/d}=(g^{n/c})^{c/d}$，故 $H_d\subseteq H_c$。反之，若包含关系成立，则 $H_d$ 的生成元作为 $H_c$ 的元素，其 $c$ 次幂为 $e$；由它的阶为 $d$ 及命题\ref{prop:order-divisibility}得 $d\mid c$。
\end{proof}

全部子群及其包含关系组成\keyterm{子群格（subgroup lattice）}：两个子群的最大公共子群是它们的交，包含两者的最小子群是由它们的并生成的子群。有限循环群的子群格由正因子的整除关系完全确定，特别地
\[
  H_d\cap H_c=H_{\gcd(d,c)},\qquad
  \langle H_d\cup H_c\rangle=H_{\operatorname{lcm}(d,c)}.
\]
这些等式分别对应最大公因数与最小公倍数的整除性质。它们在一般群中没有相同形式，不能仅凭子群阶就推断子群相等。

\begin{example}[十二阶循环群]
在加法群 $\mathbb Z/12\mathbb Z$ 中，全部子群依次为
\[
  \langle\overline0\rangle,\quad\langle\overline6\rangle,\quad
  \langle\overline4\rangle,\quad\langle\overline3\rangle,\quad
  \langle\overline2\rangle,\quad\langle\overline1\rangle,
\]
阶分别为 $1,2,3,4,6,12$。整个群的生成元是 $\overline1,\overline5,\overline7,\overline{11}$。阶为 $3$ 与阶为 $4$ 的两个子群之交为平凡子群，但它们共同生成整个群，因为 $\gcd(3,4)=1$、$\operatorname{lcm}(3,4)=12$。
\end{example}

还可以按元素的阶对 $n$ 阶循环群计数。对每个 $d\mid n$，阶为 $d$ 的元素恰为唯一 $d$ 阶子群的生成元，所以
\begin{equation}\label{eq:totient-divisor-sum}
  \sum_{d\mid n}\varphi(d)=n.
\end{equation}
这里每个元素只计入其自身阶所对应的一项。

\section{置换与循环分解}\label{sec:permutation-cycles}

\begin{definition}[对称群与循环]
集合 $\{1,\ldots,n\}$ 到自身的双射称为一个\keyterm{置换（permutation）}。全部这类置换在复合下组成\keyterm{对称群（symmetric group）}$S_n$，其中 $n\geq1$。对两两不同的 $a_1,\ldots,a_k$，记号 $(a_1\ a_2\ \cdots\ a_k)$ 表示把 $a_1$ 送到 $a_2$、依次把 $a_j$ 送到 $a_{j+1}$、把 $a_k$ 送回 $a_1$，并固定其余元素的置换，称为 $k$-循环（$k$-cycle）。
\end{definition}

$S_n$ 的阶为 $n!$，因为依次选择 $1,2,\ldots,n$ 的像时分别有 $n,n-1,\ldots,1$ 种选择。本书约定 $\sigma\tau=\sigma\circ\tau$，即先作用右侧的 $\tau$。循环写法可以循环移动起点，例如 $(1\ 3\ 2)=(3\ 2\ 1)$，但不能任意重排其中元素。

\begin{theorem}[不相交循环分解]\label{thm:disjoint-cycles}
每个 $\sigma\in S_n$ 都可分解为两两不相交的循环之积。若固定点也写为 $1$-循环，则除各循环的排列顺序及每个循环起点的选择外，分解唯一。
\end{theorem}
\begin{proof}
从一个元素 $a$ 出发考察 $a,\sigma(a),\sigma^2(a),\ldots$。集合有限，必出现重复；由 $\sigma$ 可逆，重复意味着某个正整数次迭代回到 $a$。取最小这样的次数，之前各项两两不同，并形成一个循环。若尚有未出现的元素，则从该元素重复此过程。两个所得集合若相交，应用适当的 $\sigma$ 或 $\sigma^{-1}$ 就能将其中任意元素移至交点，因此两个集合实际相同；故不同集合不相交。有限步后覆盖全部元素，所得循环的乘积在每个元素上的作用与 $\sigma$ 一致。每个元素的迭代序列由 $\sigma$ 唯一确定，也就确定了它所属的循环。
\end{proof}

不相交循环可以交换，因为一个循环移动的元素全被另一个固定。循环的逆为 $(a_1\ a_k\ a_{k-1}\ \cdots\ a_2)$。固定点通常省略，恒等置换记为 $e$。

\begin{proposition}[置换的阶]\label{prop:permutation-order}
若 $\sigma$ 的不相交循环长度为 $k_1,\ldots,k_r$，则
\[
  \operatorname{ord}(\sigma)=\operatorname{lcm}(k_1,\ldots,k_r).
\]
\end{proposition}
\begin{proof}
在一个长度为 $k_j$ 的循环上，$\sigma^m$ 固定所有元素当且仅当 $k_j\mid m$。因此 $\sigma^m=e$ 当且仅当所有 $k_j$ 都整除 $m$，最小正解正是最小公倍数。恒等置换的所有循环长度均为 $1$，阶亦为 $1$。
\end{proof}

\begin{example}[复合与元素阶]\label{ex:permutation-products}
设 $\sigma=(1\ 3\ 2)(4\ 5)$、$\tau=(1\ 4\ 2)$。按从右向左作用逐个追踪元素，得到
\[
  \sigma\tau=(1\ 5\ 4)(2\ 3),\qquad
  \tau\sigma=(1\ 3)(2\ 4\ 5).
\]
例如在第一个复合中，$1\xmapsto{\tau}4\xmapsto{\sigma}5$，然后 $5\mapsto4\mapsto1$ 形成一个三循环。两个复合不同，显示 $S_5$ 不交换。$\sigma$ 的阶为 $\operatorname{lcm}(3,2)=6$，而不是两段循环长度之和。
\end{example}

\section{奇偶性与交错群}\label{sec:permutation-sign}

长度为 $2$ 的循环称为\keyterm{对换（transposition）}。由
\begin{equation}\label{eq:cycle-transpositions}
  (a_1\ a_2\ \cdots\ a_k)
  =(a_1\ a_k)(a_1\ a_{k-1})\cdots(a_1\ a_2)
\end{equation}
及不相交循环分解，每个置换都可写成对换的乘积。该表达并不唯一，例如可在任意位置插入两个相同的对换。要定义奇偶性，必须证明所有分解的对换数量具有相同奇偶。

对 $\sigma\in S_n$，定义\keyterm{逆序数（inversion number）}
\[
  \operatorname{inv}(\sigma)
  =\bigl|\{(i,j):1\leq i<j\leq n,\ \sigma(i)>\sigma(j)\}\bigr|.
\]

\begin{lemma}[对换改变逆序数的奇偶性]\label{lem:transposition-parity}
若 $t$ 为对换，则 $\operatorname{inv}(\sigma t)$ 与 $\operatorname{inv}(\sigma)$ 的奇偶性相反。
\end{lemma}
\begin{proof}
先设 $t=(i\ i+1)$。右乘 $t$ 把序列 $\sigma(1),\ldots,\sigma(n)$ 的两个相邻位置交换。任意其他位置都在两者同一侧，它与两者形成的逆序总数不变；只有这两个位置之间的逆序状态改变，所以逆序数改变 $1$ 或 $-1$。

一般的 $t=(i\ j)$，其中 $i<j$，可写成
\[
  (i\ j)=(i\ i+1)(i+1\ i+2)\cdots(j-1\ j)
  (j-2\ j-1)\cdots(i\ i+1).
\]
这是 $2(j-i)-1$ 个相邻对换之积；对其中每一步应用前面的结论，奇偶性改变奇数次，故总效果仍是改变奇偶性。$j=i+1$ 时后半部分为空乘积。
\end{proof}

\begin{theorem}[置换奇偶性的良定义性]\label{thm:permutation-parity}
若 $\sigma=t_1\cdots t_m$ 是对换分解，则 $m\equiv\operatorname{inv}(\sigma)\pmod2$。因而分解中对换数的奇偶性只由 $\sigma$ 决定。
\end{theorem}
\begin{proof}
从恒等置换开始依次右乘 $t_1,\ldots,t_m$，由引理\ref{lem:transposition-parity}，每次使逆序数的奇偶性改变一次。恒等置换的逆序数为 $0$，结论随即得到。
\end{proof}

\begin{definition}[符号与交错群]
定义置换的\keyterm{符号（sign）}$\operatorname{sgn}(\sigma)=(-1)^{\operatorname{inv}(\sigma)}$。符号为 $1$ 的置换称为偶置换，符号为 $-1$ 的称为奇置换。全部偶置换组成的集合记为 $A_n$。
\end{definition}

\begin{proposition}[符号的乘法性]\label{prop:sign-multiplicative}
对任意 $\sigma,\tau\in S_n$，有
\begin{equation}\label{eq:sign-multiplicative}
  \operatorname{sgn}(\sigma\tau)
  =\operatorname{sgn}(\sigma)\operatorname{sgn}(\tau).
\end{equation}
因此 $A_n\leq S_n$，称为\keyterm{交错群（alternating group）}。当 $n\geq2$ 时，$|A_n|=n!/2$。
\end{proposition}
\begin{proof}
分别把 $\sigma,\tau$ 写成 $r,s$ 个对换的乘积，拼接后给出 $\sigma\tau$ 的 $r+s$ 个对换分解，由定理\ref{thm:permutation-parity}得到式\eqref{eq:sign-multiplicative}。恒等置换的符号为 $1$；由 $\sigma\sigma^{-1}=e$ 得 $\operatorname{sgn}(\sigma^{-1})=\operatorname{sgn}(\sigma)$，所以偶置换对乘法及取逆封闭，构成子群。固定一个对换 $t$，映射 $\sigma\mapsto t\sigma$ 在奇、偶置换间互换，且自身为逆，因此两类数量相同，各占 $n!/2$。
\end{proof}

式\eqref{eq:cycle-transpositions}给出 $k$-循环的符号为 $(-1)^{k-1}$。若 $\sigma$ 的不相交循环个数为 $c$，其中包含固定点对应的 $1$-循环，则
\begin{equation}\label{eq:sign-cycle-count}
  \operatorname{sgn}(\sigma)=(-1)^{n-c}.
\end{equation}
例如三循环是偶置换，四循环是奇置换。符号描述的是置换的结构，不是置换阶的奇偶性；两个不相交对换之积就是阶为 $2$ 的偶置换。

\section{例题、练习与自检}

\begin{example}[阿贝尔群未必循环]\label{ex:klein-four}
在 $S_4$ 中令 $a=(1\ 2)(3\ 4)$、$b=(1\ 3)(2\ 4)$、$c=(1\ 4)(2\ 3)$。逐点计算得到 $ab=ba=c$、$bc=cb=a$、$ca=ac=b$，并且 $a^2=b^2=c^2=e$。所以 $V_4=\{e,a,b,c\}$ 是四阶阿贝尔群，称为\keyterm{克莱因四元群（Klein four-group）}。它不循环，因为任一元素最多生成两个元素；它也不与四阶循环群同构，因为后者含四阶元素。三种非单位元都是偶置换，故 $V_4\leq A_4$。
\end{example}

\begin{enumerate}
  \item\label{exer:cyclic-power-orders} 设 $g$ 的阶为 $84$，求 $g^{30}$ 的阶。若 $g$ 无限阶，求 $\langle g^6\rangle\cap\langle g^{10}\rangle$ 以及 $\langle g^6,g^{10}\rangle$。
  \item\label{exer:cyclic-eighteen} 列出 $\mathbb Z/18\mathbb Z$ 的所有子群与整个群的生成元，并判断哪些子群之间存在包含关系。
  \item\label{exer:cyclic-involution} 证明偶数阶循环群恰有一个二阶元素，奇数阶循环群没有二阶元素。借此再次说明 $V_4$ 不循环。
  \item\label{exer:permutation-six} 设 $\sigma=(1\ 4\ 6\ 2)(3\ 5)\in S_6$，求 $\sigma^{-1}$、$\sigma^2$、$\operatorname{ord}(\sigma)$ 与 $\operatorname{sgn}(\sigma)$。
  \item\label{exer:symmetric-four-orders} 列出 $S_4$ 中元素所有可能的阶，并证明没有遗漏。能否找到一个六阶元素？
  \item\label{exer:symmetric-generators} 证明每个置换都是相邻对换 $(1\ 2),(2\ 3),\ldots,(n-1\ n)$ 的乘积，再证明当 $n\geq2$ 时 $S_n$ 由 $(1\ 2)$ 与 $(1\ 2\ \cdots\ n)$ 生成。
  \item\label{exer:small-alternating-groups} 证明 $A_3$ 是循环群，写出其全部元素；判断 $A_4$ 是否为阿贝尔群，并用具体乘积支持结论。
  \item\label{exer:totient-and-parity} 利用式\eqref{eq:totient-divisor-sum}计算 $\varphi(12)$，再用生成元判别直接核验。说明为什么某一次对换分解的长度，不能在尚未证明定理\ref{thm:permutation-parity}时就被当作置换的不变量。
\end{enumerate}

\textbf{部分提示与结果。} 第\ref{exer:cyclic-power-orders}题结果为 $14$、$\langle g^{30}\rangle$ 和 $\langle g^2\rangle$，最后一项使用贝祖等式。第\ref{exer:cyclic-eighteen}题子群的阶为 $1,2,3,6,9,18$，生成元为 $\overline1,\overline5,\overline7,\overline{11},\overline{13},\overline{17}$。第\ref{exer:permutation-six}题平方为 $(1\ 6)(2\ 4)$，阶为 $4$，符号为 $1$。第\ref{exer:symmetric-four-orders}题按不相交循环长度对 $4$ 的分拆分析，可能的阶恰为 $1,2,3,4$。第\ref{exer:symmetric-generators}题令 $c=(1\ 2\ \cdots\ n)$，计算 $c^{i-1}(1\ 2)c^{-(i-1)}$。

循环群的分类给出了“整体结构由一个生成元控制”的完整例子，置换的循环分解则把计算化为若干互不干扰的部分。下一章将对任意子群构造陪集，把这些具体计算扩展为一般的有限群计数工具。

\chapter{陪集、指数与拉格朗日定理}\label{ch:cosets}

\begin{description}
  \item[学习目标] 用陪集划分理解整除结论，区分集合的分块与商群的运算。
  \item[先修知识] 第\ref{ch:groups}、\ref{ch:cyclic-permutation}章。
\end{description}

\section{左陪集、右陪集与等价关系}
对 $H\leq G$ 及 $g\in G$ 定义左陪集 $gH$、右陪集 $Hg$。证明左陪集相交即相等，并且每个左陪集与 $H$ 等势。定义指数 $[G:H]$ 为左陪集的个数，说明左、右陪集数量相同，但相应陪集不必相等。

\section{拉格朗日定理及其边界}
对有限群 $G$ 与子群 $H$，证明拉格朗日定理（Lagrange's theorem）$|G|=[G:H]|H|$。推出元素阶整除群阶、素数阶群必为循环群。以 $A_4$ 没有 $6$ 阶子群作为反向命题失败的例子，其完整证明放在正规子群工具建立后完成。

\section{指数的乘法公式}
在 $K\leq H\leq G$ 且相关指数有限时，证明 $[G:K]=[G:H][H:K]$。通过选择陪集代表元给出可核查的计数论证。说明陪集代表元通常不唯一，计数结果却不依赖代表元的具体选择。

\section{数论应用与例题练习}
在模 $n$ 可逆剩余类构成的乘法群中推导欧拉定理，并以素数模数得到费马小定理。计算 $S_3$ 中不同子群的左右陪集，比较陪集数量与子群个数。自检要求解释为什么仅有 $|H|\mid |G|$ 不能保证某个指定子集是子群，也不能保证每个因子都对应子群。

\chapter{正规子群、商群与同态定理}\label{ch:quotient-groups}

\begin{description}
  \item[学习目标] 理解正规性控制商群乘法的良定义性，能够用核与像识别商群。
  \item[先修知识] 第\ref{ch:cosets}章。
\end{description}

\section{群同态、核与像}
在定义\ref{def:group-isomorphism}的基础上系统引入群同态，证明同态保持单位元、逆元和整数次幂。对同态 $f:G\to H$ 定义核 $\ker f$ 与像 $\operatorname{im}f$，证明单射等价于核只含单位元。结合循环群、符号映射和行列式映射计算核与像。

\section{正规子群与商群构造}
定义正规子群（normal subgroup）$N\trianglelefteq G$，证明共轭不变与左右陪集相等的等价性。证明陪集乘法 $(gN)(hN)=ghN$ 良定义当且仅当 $N$ 正规，再定义商群（quotient group）$G/N$。说明正规性不同于交换性，并证明任意同态的核正规。

\section{商映射的分解性质与第一同构定理}
对正规子群 $N\trianglelefteq G$，定义自然投影 $\pi:G\to G/N$。若同态 $f:G\to H$ 满足 $N\subseteq\ker f$，证明存在唯一同态 $\overline f:G/N\to H$ 使 $f=\overline f\circ\pi$。由此证明第一同构定理 $G/\ker f\cong\operatorname{im}f$，并强调需构造映射、验证良定义，再证明双射。

\section{其余同构定理与对应定理}
对 $H\leq G$、$N\trianglelefteq G$，证明 $HN/N\cong H/(H\cap N)$。对 $N\subseteq K$ 且 $N,K\trianglelefteq G$，证明 $(G/N)/(K/N)\cong G/K$。建立 $G/N$ 的子群与 $G$ 中包含 $N$ 的子群之间的对应，同时说明正规性如何传递。

\section{例题、练习与自检}
计算 $\mathbb Z/m\mathbb Z$、$S_n/A_n$（$n\geq2$）及若干矩阵群同态的像。证明指数为 $2$ 的子群必正规，完成 $A_4$ 不存在 $6$ 阶子群的反例。安排一个非正规子群的陪集乘法失败算例，使良定义性成为实际检查步骤。

\chapter{群作用与计数}\label{ch:actions}

\begin{description}
  \item[学习目标] 把群元素看作集合上的变换，以轨道与稳定子组织计数和结构证明。
  \item[先修知识] 第\ref{ch:quotient-groups}章。
\end{description}

\section{群作用、轨道与稳定子}
定义群 $G$ 在集合 $X$ 上的左作用（group action），以及忠实作用、传递作用。对 $x\in X$ 定义轨道 $Gx$ 与稳定子 $G_x$，证明轨道划分 $X$，并构造轨道与左陪集集合之间的双射；当 $G$ 有限时得到轨道--稳定子公式 $|Gx|=[G:G_x]$。

\section{置换表示与凯莱定理}
说明一个群作用等价于一个从 $G$ 到 $X$ 的全体置换所成群的同态。利用群对自身的左乘作用证明凯莱定理（Cayley's theorem）：每个群都同构于某个置换群的子群。比较左乘作用、对子群陪集的作用与几何对称作用，分别计算其核。

\section{共轭作用与类方程}
定义共轭类（conjugacy class）、中心 $Z(G)$ 和元素的中心化子 $C_G(g)$。对有限群写出类方程，并由共轭作用解释共轭类大小等于 $[G:C_G(g)]$。证明非平凡有限 $p$ 群的中心非平凡，其中 $p$ 为素数、$p$ 群指阶为 $p$ 的幂的群；进一步证明 $p^2$ 阶群必为阿贝尔群。

\section{伯恩赛德引理与计数应用}
对有限群 $G$ 作用于有限集合 $X$，令 $X^g=\{x\in X:g\cdot x=x\}$，以双重计数证明伯恩赛德引理（Burnside's lemma）：轨道数为 $|G|^{-1}\sum_{g\in G}|X^g|$。用正方形顶点染色说明“按旋转视为相同”和“按全部对称视为相同”对应不同的作用群。

\section{例题、练习与自检}
计算 $S_3$ 的共轭类与类方程，求正方形顶点二色染色在不同等价约定下的数量。练习要求从文字描述中写出作用集合与作用群，判断作用是否忠实或传递，并通过稳定子计算轨道，避免直接把对象总数除以群阶。

\chapter{西罗定理与有限群}\label{ch:sylow}

\begin{description}
  \item[学习目标] 利用素数幂阶子群的存在、共轭和数量限制研究有限群。
  \item[先修知识] 第\ref{ch:actions}章，尤其是轨道计数与共轭作用。
\end{description}

\section{柯西定理与素数幂子群}
证明柯西定理（Cauchy's theorem）：若素数 $p$ 整除有限群 $G$ 的阶，则 $G$ 含有一个 $p$ 阶元素。用循环移位作用给出一种证明，明确作用集合及其不动点。比较它与拉格朗日定理，说明其结论只处理素数阶而非任意因子阶。

\section{西罗定理的三个结论}
设 $|G|=p^a m$，其中 $a\geq1$ 且 $p\nmid m$。定义西罗 $p$ 子群（Sylow $p$-subgroup）为 $p^a$ 阶子群。完整陈述并证明：西罗 $p$ 子群存在；每个 $p$ 子群包含于某个西罗 $p$ 子群且所有西罗 $p$ 子群互为共轭；其数目 $n_p$ 满足 $n_p\mid m$ 与 $n_p\equiv1\pmod p$。

\section{正规化子与证明的组织}
对 $P\leq G$ 定义正规化子 $N_G(P)=\{g\in G:gPg^{-1}=P\}$。在共轭作用下推出 $n_p=[G:N_G(P)]$，并证明西罗 $p$ 子群正规当且仅当 $n_p=1$。将存在性、包含性、共轭性和数量公式分步证明，标明每一步使用的作用与计数条件。

\section{小阶群的分析与方法边界}
分析 $6$、$10$、$15$ 阶群的西罗子群数量，以 $15$ 阶群必循环说明数量限制、正规性与直积思想的结合。对 $12$ 阶群展示仅凭同余与整除条件可能保留多个候选值；进一步结论需要元素计数、置换作用或结构分析，不能把候选值当成已实现的群。

\section{例题、练习与自检}
安排判断给定阶数的群是否必有非平凡真正规子群、计算 $S_4$ 的西罗子群以及排除某些单群阶数的题目。自检要求能够说明“存在某个正规子群”与“完整分类全部群”之间的区别，并检查每次使用 $n_p\mid m$ 时 $m$ 的取值。

\chapter{直积与有限阿贝尔群}\label{ch:abelian}

\begin{description}
  \item[学习目标] 理解群的分解，掌握有限阿贝尔群分类及其同构判别。
  \item[先修知识] 第\ref{ch:quotient-groups}章；有限群相关例子结合第\ref{ch:sylow}章。
\end{description}

\section{外直积与内直积}
定义直积（direct product）$G\times H$ 的逐分量运算。对群 $G$ 的两个正规子群 $H,K$，在 $G=HK$ 且 $H\cap K=\{e\}$ 时证明 $G\cong H\times K$；先证明两子群的元素彼此交换，再验证乘法映射为同构。推广到有限个因子的情形。

\section{循环群的直积分解}
证明 $\mathbb Z/m\mathbb Z\times\mathbb Z/n\mathbb Z$ 为循环群当且仅当 $\gcd(m,n)=1$。用素数幂因子分解有限循环群，比较 $\mathbb Z/12\mathbb Z$ 与 $\mathbb Z/6\mathbb Z\times\mathbb Z/2\mathbb Z$ 的元素阶。将这类分解与第\ref{ch:ideals}章的环版本中国剩余定理衔接。

\section{有限阿贝尔群结构定理}
陈述并证明有限阿贝尔群是有限个素数幂阶循环群的直积，且这些因子除排列顺序外唯一；再改写为非平凡循环因子的阶逐项整除的形式。分别引入初等因子（elementary divisors）和不变因子（invariant factors）。证明路线依次为素数分量分解、阿贝尔 $p$ 群的循环直因子分解、由各 $p^k$ 阶挠子群大小恢复因子重数，不依赖后面的模论。

\section{半直积初步（选学）}
定义自同构群 $\operatorname{Aut}(N)$。给定群同态 $\varphi:H\to\operatorname{Aut}(N)$，在 $N\times H$ 上用 $(n,h)(n',h')=(n\varphi(h)(n'),hh')$ 定义半直积（semidirect product）$N\rtimes_\varphi H$，验证群公理。用二面体群说明正规因子与补子群的元素不必彼此交换，直积只对应作用平凡的特例。

\section{例题、练习与自检}
列出所有 $72$ 阶阿贝尔群的同构类型，转换初等因子与不变因子表示，并用元素阶或挠子群大小区分同阶群。比较 $S_3$ 的半直积描述与 $\mathbb Z/6\mathbb Z$ 的直积描述；强调有限阿贝尔群分类不适用于一般有限群。

\chapter{群列、单群与可解群}\label{ch:solvable-groups}

\begin{description}
  \item[学习目标] 用逐层取商刻画群的复杂性，为根式可解性准备群论语言。
  \item[先修知识] 第\ref{ch:quotient-groups}、\ref{ch:actions}章；例题使用西罗定理。
\end{description}

\section{群列与单群}
约定次正规列（subnormal series）为有限链 $G=G_0\supseteq G_1\supseteq\cdots\supseteq G_r=\{e\}$，其中每个 $G_{i+1}\trianglelefteq G_i$；若每项都正规于 $G$，则称为正规列。定义因子群 $G_i/G_{i+1}$ 与单群（simple group），要求单群非平凡且没有非平凡真正规子群。用素数阶循环群建立最初例子。

\section{换位子与导出列}
固定换位子（commutator）$[a,b]=aba^{-1}b^{-1}$ 的约定，定义导出子群 $G'=[G,G]$ 为所有换位子生成的子群。证明 $G'$ 正规且 $G/G'$ 阿贝尔，并刻画其最小性：对 $N\trianglelefteq G$，商群 $G/N$ 阿贝尔当且仅当 $G'\subseteq N$。递归定义 $G^{(0)}=G$、$G^{(i+1)}=(G^{(i)})'$。

\section{可解群及其封闭性质}
定义可解群（solvable group）为导出列在有限步后达到平凡群的群，证明它等价于存在阿贝尔因子的次正规列。证明可解性对子群、商群及扩张封闭；扩张性质具体指对 $N\trianglelefteq G$，若 $N$ 与 $G/N$ 均可解，则 $G$ 可解。区分阿贝尔群、可解群与一般群。

\section{交错群与不可解性的第一个例子}
计算 $S_3$、$S_4$ 的导出列或构造阿贝尔因子的群列。通过 $A_5$ 的共轭类大小证明其为非阿贝尔单群，再推出 $A_5$ 与 $S_5$ 不可解。一般的 $A_n$（$n\geq5$）单性证明列为选学，并说明它在高次方程问题中的作用。

\section{合成列与练习规划}
定义合成列（composition series），将有限群合成列的存在性和若尔当--赫尔德定理（Jordan--Hölder theorem）列为选学：合成因子的同构类型及重数与合成列选择无关。核心练习要求计算换位子、检验一条群列是否满足条件，并给出可解但非阿贝尔的例子。

\part{环与多项式}

\chapter{环、整环、域与环同态}\label{ch:rings}

\begin{description}
  \item[学习目标] 区分环、整环与域，将群同态的思想推广到同时保持两种运算的映射。
  \item[先修知识] 第\ref{ch:quotient-groups}章。
\end{description}

\section{环的公理与基本例子}
由整数、多项式、矩阵和函数的逐点运算引出环公理，回顾环的加法必须组成阿贝尔群。区分交换环与非交换环，检查子环是否包含原环的单位元。以 $\mathbb Z/n\mathbb Z$ 的剩余类运算说明同一集合可同时承载加法群和环结构。

\section{单位、零因子与整环}
定义单位（unit）、单位群 $R^\times$、左零因子和右零因子；在交换环中统一讨论零因子。定义整环（integral domain）为非零、交换且没有非零零因子的环，证明整环中的消去律。以矩阵环、整数环和模合数剩余类环区分可逆性与非零性。

\section{域、特征与素域}
定义域为非零且每个非零元素可逆的交换环。定义环的特征（characteristic）$\operatorname{char}R$ 为使 $n1_R=0$ 的最小正整数 $n$，不存在时记为 $0$。证明整环的特征为 $0$ 或素数，有限整环必为域。对素数 $p$ 记 $\mathbb F_p=\mathbb Z/p\mathbb Z$，并证明任意域的最小子域同构于 $\mathbb Q$ 或某个 $\mathbb F_p$。

\section{环同态与结构保持}
定义保持加法、乘法与单位元的环同态及环同构。比较自然剩余类映射、包含映射和多项式求值映射。说明环同态的核对加法形成子群，但一般不是本书约定下的含幺子环；这正是引入理想的动机。

\section{例题、练习与自检}
计算 $\mathbb Z/12\mathbb Z$ 的单位与非零零因子，验证高斯整数环 $\mathbb Z[i]=\{a+bi:a,b\in\mathbb Z\}$ 为整环。练习要求给出非交换环、非域整环及非整环交换环，并检查给定映射是否满足全部环同态条件。

\chapter{理想、商环与中国剩余定理}\label{ch:ideals}

\begin{description}
  \item[学习目标] 理解理想控制商环运算，能够用极大理想构造域并分解适当的商环。
  \item[先修知识] 第\ref{ch:rings}章及群同构定理。
\end{description}

\section{理想与生成理想}
定义左理想、右理想和双边理想（ideal）；构造商环时均使用双边理想。说明理想是加法子群并满足乘法吸收条件，不能与子环混为一谈。对交换环定义主理想 $(a)$、有限生成理想 $(a_1,\ldots,a_r)$，介绍理想的和、交与积。

\section{商环及环同构定理}
对环 $R$ 的双边理想 $I$，验证 $R/I$ 的加法与乘法良定义。证明环同态的核为双边理想，以及 $R/\ker f\cong\operatorname{im}f$。给出对应定理和迭代取商公式，并用求值映射展示“关系式被设为零”的商环解释。

\section{素理想与极大理想}
本节设 $R$ 为交换环，且所讨论理想均为真理想。定义素理想（prime ideal）与极大理想（maximal ideal），证明 $P$ 为素理想当且仅当 $R/P$ 为整环，$M$ 为极大理想当且仅当 $R/M$ 为域。推出极大理想必为素理想，并用 $\mathbb Z$ 的零理想说明逆命题不成立。

\section{中国剩余定理}
对交换环 $R$ 中两两互素的有限个理想 $I_1,\ldots,I_r$，其中互素指 $I_i+I_j=R$（$i\ne j$），证明自然同态满射且其核为各理想之交，并证明此时交等于积，从而得到
\begin{equation}\label{eq:ring-crt}
  R/(I_1\cdots I_r)\cong\prod_{j=1}^{r}R/I_j.
\end{equation}
先处理两个理想，再归纳推广；对整数模互素数的情形，给出由贝祖等式构造原像的方法。

\section{例题、练习与自检}
解一组模 $3$、$5$、$7$ 的联立同余，分析 $\mathbb Z/60\mathbb Z$ 的直积分解。比较 $\mathbb R[x]/(x^2+1)$ 与 $\mathbb R[x]/(x^2-1)$，说明商环分别表现为域与含零因子的环。练习要求写出同构的具体映射，并指出理想不互素时式\eqref{eq:ring-crt}的满射性为何可能失败。

\chapter{整除性与唯一分解}\label{ch:factorization}

\begin{description}
  \item[学习目标] 识别整数分解论中哪些性质能够推广到整环，区分不同强度的分解条件。
  \item[先修知识] 第\ref{ch:ideals}章。
\end{description}

\section{整除、相伴、不可约元与素元}
本章设 $R$ 为整环。定义整除关系、相伴元（associates）、不可约元（irreducible element）与素元（prime element），两类特殊元素均要求非零且不是单位。证明素元必不可约，说明一般整环中逆命题不成立。最大公因数若存在，仅在相伴意义下唯一。

\section{欧几里得整环与主理想整环}
定义欧几里得整环（Euclidean domain）及其在非零元素上的取值于非负整数的除法函数，要求每次带余除法的非零余式具有更小函数值。定义主理想整环（principal ideal domain, PID），证明欧几里得整环必为主理想整环，并由理想 $(a,b)$ 推导最大公因数的贝祖表示。

\section{唯一分解整环及其判别}
定义唯一分解整环（unique factorization domain, UFD）：每个非零非单位均能分解为有限个不可约元，分解除顺序与相伴替换外唯一。证明主理想整环必为唯一分解整环，分别处理分解的存在与唯一性；总结“欧几里得整环 $\Rightarrow$ 主理想整环 $\Rightarrow$ 唯一分解整环”，避免未经证明倒用箭头。

\section{分式域与典型反例}
对整环 $R$，用分母非零的有序对及交叉相乘等价关系构造分式域（field of fractions）$\operatorname{Frac}(R)$。以高斯整数中的范数说明欧几里得算法，以 $\mathbb Z[\sqrt{-5}]$ 中 $6=2\cdot3=(1+\sqrt{-5})(1-\sqrt{-5})$ 分析唯一分解失效；必须验证各因子不可约且不相伴，而不只展示两种乘积。

\section{例题、练习与自检}
在 $\mathbb Z[i]$ 中做带余除法并求最大公因数，练习从生成理想转化为整除关系。将 $F[x,y]$（$F$ 为域）作为唯一分解整环不必为主理想整环的例子，其唯一分解性待第\ref{ch:polynomials}章建立多项式分解定理后补证；本章先分析理想 $(x,y)$ 为什么不是主理想。

\chapter{多项式环与不可约性}\label{ch:polynomials}

\begin{description}
  \item[学习目标] 掌握域上多项式的除法、分解和不可约性判别，为添加方程的根建立代数基础。
  \item[先修知识] 第\ref{ch:factorization}章，初等多项式运算。
\end{description}

\section{形式多项式、次数与求值}
对交换环 $R$ 定义形式多项式环 $R[x]$，说明形式多项式与其诱导的多项式函数不同；在有限域上不同多项式可能给出同一函数。对整环证明非零多项式乘积的次数可加，并构造求值同态。定义有理函数域 $F(t)=\operatorname{Frac}(F[t])$，其中 $F$ 为域、$t$ 为不定元。

\section{域上多项式的带余除法}
设 $F$ 为域，证明对 $f,g\in F[x]$ 且 $g\ne0$，存在唯一 $q,r\in F[x]$ 使 $f=qg+r$，其中 $r=0$ 或 $\deg r<\deg g$。由此证明 $F[x]$ 为欧几里得整环，建立多项式欧几里得算法、贝祖等式、余式定理与因式定理。

\section{高斯引理与不可约性判别}
在唯一分解整环 $R$ 中定义多项式的内容与本原多项式（primitive polynomial），证明高斯引理（Gauss's lemma），进而证明 $R[x]$ 仍为唯一分解整环。对于正次数的本原多项式，比较它在 $R[x]$ 与 $\operatorname{Frac}(R)[x]$ 中的不可约性。系统安排有理根判别、约化模素数法和艾森斯坦判别法（Eisenstein's criterion），完整说明首项系数与常数项的限制。

\section{根、重根与商环构造域}
证明域上非零 $d$ 次多项式至多有 $d$ 个不同根，定义形式导数与根的重数。对非恒定多项式 $f\in F[x]$，证明它在分裂域中无重根当且仅当 $\gcd(f,f')=1$；分裂域的构造在第\ref{ch:splitting}章完成。证明 $F[x]/(f)$ 为域当且仅当 $f$ 不可约，并说明商环中 $x$ 的剩余类是 $f$ 的一个根。

\section{例题、练习与自检}
判断 $x^4+1$ 在 $\mathbb Q[x]$ 中的不可约性，构造 $\mathbb F_2[x]/(x^2+x+1)$ 并计算元素的逆。比较同一多项式在 $\mathbb Q$、$\mathbb R$ 与 $\mathbb C$ 上的分解。练习中特别安排“没有根但仍可约”的四次多项式，说明无根判别只直接适用于二次、三次多项式。

\part{域与伽罗瓦理论}

\chapter{域扩张与代数元素}\label{ch:extensions}

\begin{description}
  \item[学习目标] 把添加元素转化为域与向量空间的构造，使用最小多项式计算扩张次数。
  \item[先修知识] 第\ref{ch:polynomials}章，以及线性空间的基与维数。
\end{description}

\section{域扩张、生成域与扩张次数}
定义域扩张（field extension）$K/F$ 表示 $F$ 是 $K$ 的子域，将 $K$ 看作 $F$ 上的向量空间，并记 $[K:F]=\dim_FK$。定义生成域 $F(S)$ 与单扩张 $F(\alpha)$，区分它们与只允许有限次加法、乘法得到的子环 $F[S]$、$F[\alpha]$。

\section{代数元素与最小多项式}
定义代数元素（algebraic element）、超越元素（transcendental element）和首一最小多项式 $m_{\alpha,F}(x)$。用求值映射的核证明最小多项式存在时唯一且不可约，建立
\begin{equation}\label{eq:simple-extension}
  F(\alpha)=F[\alpha]\cong F[x]/(m_{\alpha,F}),
  \qquad [F(\alpha):F]=\deg m_{\alpha,F}.
\end{equation}
说明式\eqref{eq:simple-extension}的前提是 $\alpha$ 在 $F$ 上代数；超越情形对应有理函数域。

\section{塔式公式与基的构造}
对域塔 $F\subseteq K\subseteq L$，在 $[K:F]$ 和 $[L:K]$ 有限时，通过两组基的乘积构造证明 $[L:F]=[L:K][K:F]$。计算二次扩张与连续添加根的扩张，解释最小多项式会随基域改变，不能直接把每个已知多项式的次数相乘当作扩张次数。

\section{有限扩张与代数扩张}
定义代数扩张，证明有限扩张必代数、有限个代数元素生成有限扩张，以及代数性的传递性。区分“每个元素分别满足某个多项式”与“整个扩张具有有限维数”；代数而非有限的例子在第\ref{ch:splitting}章的代数闭包讨论中补充。

\section{例题、练习与自检}
计算 $\mathbb Q(\sqrt2,\sqrt3)$ 的基与次数，求 $\sqrt2+\sqrt3$ 的最小多项式，并计算扩域中元素的逆。练习要求先给出一个消去元素的多项式，再单独证明其不可约性；不能把任何以该元素为根的多项式都称为最小多项式。

\chapter{尺规作图与二次扩张}\label{ch:constructions}

\begin{description}
  \item[学习目标] 将几何作图步骤转化为域运算，理解代数条件如何给出不可作图证明。
  \item[先修知识] 第\ref{ch:extensions}章，平面解析几何。本章为应用专题，可独立选学。
\end{description}

\section{可作点与可作数}
在欧氏平面中固定原点与单位长度，明确无刻度直尺和圆规允许的操作。定义可作点、可作实数（constructible real number），说明坐标系与初始数据的作用。先建立加、减、乘、除及对正数开平方的几何实现。

\section{交点方程与二次扩张塔}
分析直线与直线、直线与圆、圆与圆的交点方程，证明一次作图至多添加二次方程的根。建立实数 $\alpha$ 从有理初始数据可作，当且仅当它属于一条有限实域塔的顶层，其中起点为 $\mathbb Q$、每步扩张次数为 $1$ 或 $2$。分别证明必要性与充分性。

\section{次数障碍与经典问题}
由塔式公式推出可作代数数 $\alpha$ 的次数 $[\mathbb Q(\alpha):\mathbb Q]$ 必为 $2$ 的幂，并明确此条件本身不充分。用 $\sqrt[3]{2}$ 证明倍立方不可作，用三倍角关系与不可约三次多项式证明一般三等分角不可作；某些特定角仍可三等分。

\section{正多边形与论证边界}
将正多边形作图与单位根联系，正五边形作为核心例题。一般正多边形的可作判别在圆分扩张之后列为选学回访。化圆为方的讨论明确依赖 $\pi$ 的超越性；该超越性作为外部定理说明，不声称可由本章初等域论证明。

\section{例题、练习与自检}
构造若干含平方根的长度，判断给定代数数是否能由次数障碍排除。练习要求区分“不能由当前必要条件判定”与“已经证明可作”，并解释为什么证明某个角不可三等分足以否定普遍作图程序，却不能否定所有特殊角。

\chapter{分裂域、可分扩张与正规扩张}\label{ch:splitting}

\begin{description}
  \item[学习目标] 理解包含全部根、根无重复和嵌入数量之间的联系，准确掌握伽罗瓦扩张的条件。
  \item[先修知识] 第\ref{ch:extensions}章、多项式不可约性与形式导数。
\end{description}

\section{分裂域的构造与唯一性}
对域 $F$ 上的非恒定多项式 $f$，定义分裂域（splitting field）为 $f$ 分裂且由其全部根生成的扩域。通过逐次取不可约因子并添加根证明存在性，通过扩张嵌入证明分裂域在 $F$-同构意义下唯一；唯一性不表示存在唯一一个同构。

\section{可分多项式与可分扩张}
定义可分多项式（separable polynomial）为在分裂域中没有重根的多项式，代数元素可分指其最小多项式可分。代数扩张可分指其所有元素可分。证明特征为 $0$ 的域上不可约多项式可分，并用 $\mathbb F_p(t)$ 上的 $x^p-t$ 展示正特征下不可分现象，分别验证不可约性与重根性质。

\section{正规扩张与伽罗瓦扩张}
定义代数扩张 $K/F$ 正规（normal），是指 $F[x]$ 中每个在 $K$ 有根的不可约多项式都在 $K$ 完全分裂。证明有限扩张正规当且仅当它是某个多项式的分裂域。定义伽罗瓦扩张（Galois extension）为正规且可分的代数扩张；本书随后主要处理有限情形。

\section{嵌入数量与本原元}
定义固定基域 $F$ 的域嵌入，证明有限扩张 $K/F$ 到一个代数闭包中的 $F$-嵌入数不超过 $[K:F]$，等号成立当且仅当扩张可分。证明本原元定理（primitive element theorem）：有限可分扩张可由单个元素生成；说明这是一条带条件的存在性结论，不能推广为任意有限扩张都为单扩张。

\section{代数闭包与例题练习}
定义代数闭包（algebraic closure）；其一般存在性与唯一性证明涉及额外集合论工具，列为选学并注明所用原则。核心例题比较 $\mathbb Q(\sqrt[3]{2})/\mathbb Q$ 与 $x^3-2$ 的分裂域。以 $\mathbb Q$ 在 $\mathbb C$ 中的代数闭包为代数而非有限的例子，通过任意高次数的艾森斯坦多项式说明次数无界。

\chapter{有限域}\label{ch:finite-fields}

\begin{description}
  \item[学习目标] 掌握有限域的存在、唯一性、子域和自同构，能够实际构造并计算有限域。
  \item[先修知识] 第\ref{ch:splitting}章与循环群基础。
\end{description}

\section{有限域的阶与构造}
若有限域特征为素数 $p$，将其视作 $\mathbb F_p$ 上的有限维向量空间，证明其阶为 $p^n$，其中 $n\geq1$。反过来，在 $x^{p^n}-x$ 的分裂域中证明全部根组成一个含 $p^n$ 个元素的域，并证明此阶的域在同构意义下唯一，记为 $\mathbb F_{p^n}$。

\section{有限域乘法群与本原元素}
证明任意域的有限乘法子群为循环群，推出 $\mathbb F_q^\times$ 为 $q-1$ 阶循环群，其中 $q$ 为素数幂。定义有限域的本原元素（primitive element）为其乘法群生成元，并与第\ref{ch:splitting}章“生成整个扩张的本原元”区分；有限域中前者能生成域扩张，后者未必生成乘法群。

\section{弗罗贝尼乌斯自同构与子域}
在 $\mathbb F_{p^n}$ 上定义弗罗贝尼乌斯自同构（Frobenius automorphism）$\sigma(a)=a^p$，证明其阶为 $n$。证明 $\mathbb F_{p^n}$ 含有 $p^d$ 阶子域当且仅当 $d\mid n$，且对每个这样的 $d$ 子域唯一，等于 $x^{p^d}-x$ 的根集。以此显式计算全部固定 $\mathbb F_p$ 的自同构。

\section{不可约多项式与有限域运算}
对首一不可约多项式 $f\in\mathbb F_p[x]$、$\deg f=n$，把 $\mathbb F_p[x]/(f)$ 的元素表示成次数小于 $n$ 的剩余类代表元。给出加法、乘法与利用扩展欧几里得算法求逆的步骤。说明首一不可约多项式 $g\in\mathbb F_p[x]$ 整除 $x^{p^n}-x$ 当且仅当 $\deg g\mid n$。

\section{例题、练习与自检}
用不可约三次多项式构造 $\mathbb F_8$，列出其非零元素的幂，计算乘法逆元。列出 $\mathbb F_{p^{12}}$ 的全部子域并判断包含关系。练习中特别区分 $\mathbb F_{p^n}$ 与 $\mathbb Z/p^n\mathbb Z$：当 $n>1$ 时后者不是域。

\chapter{伽罗瓦对应与伽罗瓦群计算}\label{ch:galois}

\begin{description}
  \item[学习目标] 在子群与中间域之间双向转换，并从多项式根的置换计算具体扩张的对称性。
  \item[先修知识] 第\ref{ch:quotient-groups}、\ref{ch:actions}、\ref{ch:splitting}章；有限域提供一类完整例子。
\end{description}

\section{域自同构、固定域与伽罗瓦群}
对扩张 $L/F$ 定义固定 $F$ 的自同构群 $\operatorname{Aut}_F(L)$；在有限伽罗瓦情形记为 $\operatorname{Gal}(L/F)$。对自同构子群 $H$ 定义固定域 $L^H=\{a\in L:\sigma(a)=a\text{ 对所有 }\sigma\in H\}$。从自同构对根的作用说明何时可把伽罗瓦群忠实嵌入对称群。

\section{阿廷固定域定理与基本定理的准备}
证明不同域嵌入的线性无关性，作为阿廷固定域定理（Artin's fixed-field theorem）的关键引理。若 $H$ 是域 $L$ 的一个有限自同构群，证明 $L/L^H$ 为有限伽罗瓦扩张且 $[L:L^H]=|H|$。与嵌入计数结合，解释有限伽罗瓦扩张的自同构数为何恰好等于扩张次数。

\section{有限伽罗瓦理论基本定理}
设 $L/F$ 为有限伽罗瓦扩张，$G=\operatorname{Gal}(L/F)$。建立中间域 $F\subseteq E\subseteq L$ 与子群 $H\leq G$ 之间的互逆对应
\begin{equation}\label{eq:galois-correspondence}
  E\longmapsto\operatorname{Gal}(L/E),
  \qquad H\longmapsto L^H.
\end{equation}
证明对应反转包含关系；当 $E=L^H$ 时，$[L:E]=|H|$、$[E:F]=[G:H]$。进一步证明 $E/F$ 为伽罗瓦扩张当且仅当 $H\trianglelefteq G$，此时限制映射诱导 $\operatorname{Gal}(E/F)\cong G/H$。区分总是伽罗瓦的 $L/E$ 与未必伽罗瓦的 $E/F$。

\section{伽罗瓦群的实际计算}
按“分裂域、扩张次数、根的可能像、关系式限制、子群判定”的次序组织计算。完整处理双二次扩张、$x^3-2$ 的分裂域与有限域扩张。对特征为 $0$ 的域上的不可约三次多项式引入判别式（discriminant），证明判别式是否为基域中的平方分别对应 $A_3$ 与 $S_3$。

\section{例题、练习与自检}
列出 $\mathbb Q(\sqrt2,\sqrt3)/\mathbb Q$ 的所有中间域，计算 $x^3-2$ 分裂域中各子群的固定域。练习要求验证某个根的置换确实保持全部代数关系，防止把任意置换都误判为自同构；另用非正规扩张说明不能直接套用式\eqref{eq:galois-correspondence}的全部结论。

\chapter{根式扩张与方程的可解性}\label{ch:radicals}

\begin{description}
  \item[学习目标] 理解根式表达与可解群之间的联系，准确陈述一般高次方程不可用根式求解的含义。
  \item[先修知识] 第\ref{ch:solvable-groups}、\ref{ch:galois}章。本章基域均假设特征为 $0$。
\end{description}

\section{根式扩张与根式可解的定义}
定义根式扩张（radical extension）为可由有限域塔得到的扩张，其中每一步形如 $F_i=F_{i-1}(\alpha_i)$，并且存在正整数 $m_i$ 使 $\alpha_i^{m_i}\in F_{i-1}$。定义多项式根式可解，是指其全部根包含于某个根式扩张中；说明分裂域本身不必就是所选的根式扩张。

\section{单位根与圆分扩张}
定义本原 $n$ 次单位根 $\zeta_n$ 和圆分多项式（cyclotomic polynomial）$\Phi_n(x)$。证明 $\Phi_n\in\mathbb Z[x]$ 且在 $\mathbb Q[x]$ 中不可约，从而建立 $\operatorname{Gal}(\mathbb Q(\zeta_n)/\mathbb Q)\cong(\mathbb Z/n\mathbb Z)^\times$。解释添加单位根为何能控制纯根式方程的自同构。

\section{根式可解与可解群的等价}
设 $F$ 为特征 $0$ 的域，$f\in F[x]$ 为非恒定多项式，$L$ 为其分裂域。证明 $f$ 在 $F$ 上根式可解当且仅当 $\operatorname{Gal}(L/F)$ 为可解群。分别安排两个方向：对根式塔添加必要单位根并控制其正规闭包；对可解群逐层取固定域，并利用含单位根时的素数阶循环扩张构造根式。后一个步骤需单独证明，不能仅从“商群阿贝尔”直接跳到“能开根”。

\section{低次方程与一般五次方程}
结合 $S_2,S_3,S_4$ 的可解性解释二、三、四次方程存在根式解；用 $S_n$（$n\geq5$）的不可解性分析一般高次方程。若表述“一般 $n$ 次方程”，明确以代数独立系数上的有理函数域为基域，并证明其伽罗瓦群为 $S_n$。同时安排一个伽罗瓦群为 $S_5$ 的具体有理系数五次多项式，完整验证不可约性与群的判定。

\section{例题、练习与自检}
比较 $x^5-2$ 与不可根式求解的五次多项式，说明次数相同不决定根式可解性。回访尺规作图：以根式塔中的二次步骤解释作图限制，将正多边形判别列为选学。练习要求区分无通用根式公式、某个方程无根式解、方程没有根及无法数值近似这四个不同命题。

\part{模与线性代数的统一}

\chapter{模、子模与模同态}\label{ch:modules}

\begin{description}
  \item[学习目标] 用同一语言描述阿贝尔群与带线性变换的向量空间，理解域上的线性代数为什么不能直接照搬到一般环。
  \item[先修知识] 第\ref{ch:ideals}章及线性代数；与域论主线可独立衔接。
\end{description}

\section{从向量空间到模}
对环 $R$ 定义含幺左 $R$-模，明确标量作用的结合性、两种分配律和单位元条件。证明阿贝尔群与 $\mathbb Z$-模的对应，域上的模就是向量空间。对域 $F$ 上的向量空间 $V$ 及线性变换 $T:V\to V$，用 $f(x)\cdot v=f(T)v$ 定义 $F[x]$-模结构。

\section{子模、商模与同态定理}
定义子模（submodule）、模同态（module homomorphism）与商模（quotient module），建立核、像及模版本的同构定理。通过阿贝尔群与线性空间的例子比较三种同构定理的共同结构，并解释为何这里的商结构依靠加法子群对标量作用封闭。

\section{生成、自由模与基}
定义有限生成模、线性无关、基和自由模（free module）。用 $\mathbb Z/n\mathbb Z$（$n\geq2$）说明有限生成模未必自由。对一般环区分“由 $r$ 个元素生成”与“具有 $r$ 元基”；涉及基的基数唯一性时限定到本书后续使用的交换整环，不无条件推广向量空间结论。

\section{直和、挠元与正合列初步}
定义模的有限直和（direct sum），并解释无限直和要求元素只有有限个非零分量，与无限直积不同。对整环 $R$ 上的模定义挠元（torsion element）与挠子模，证明挠元全体构成子模。把短正合列（short exact sequence）作为选学语言引入，说明像等于核的含义及分裂时的直和解释。

\section{例题、练习与自检}
把 $\mathbb Z^2$ 的给定子群写成子模并计算商模，比较 $\mathbb Z$-模与 $\mathbb Q$-向量空间的生成问题。把具体矩阵转化为 $F[x]$-模，解释不变子空间就是子模。安排“无挠不一定自由”的一般反例，并明确主理想整环上的有限生成无挠模具有更强结论。

\chapter{主理想整环上的有限生成模与标准形}\label{ch:pid-modules}

\begin{description}
  \item[学习目标] 掌握有限生成模结构定理，将有限生成阿贝尔群和矩阵标准形作为同一个分解原理的应用。
  \item[先修知识] 第\ref{ch:factorization}、\ref{ch:polynomials}、\ref{ch:modules}章，以及矩阵相似与特征多项式。
\end{description}

\section{有限生成模的表示与史密斯标准形}
设 $R$ 为主理想整环，先证明有限秩自由模的子模仍为有限秩自由模，再用生成元与关系给出有限生成模的矩阵表示。通过可逆行、列变换建立史密斯标准形（Smith normal form），要求非零对角元逐项整除。说明这类矩阵等价变换不同于研究线性变换时的相似变换；在 $\mathbb Z$ 与 $F[x]$ 上给出欧几里得算法驱动的计算过程。

\section{有限生成模结构定理}
设 $M$ 为主理想整环 $R$ 上的有限生成模，证明存在 $r\geq0$ 及非零非单位 $d_1,\ldots,d_s$，满足 $d_1\mid\cdots\mid d_s$，使
\begin{equation}\label{eq:pid-module-structure}
  M\cong R^r\oplus R/(d_1)\oplus\cdots\oplus R/(d_s).
\end{equation}
允许 $s=0$；证明自由秩 $r$ 唯一，非平凡循环因子的不变因子在相伴意义下唯一。再转化为素元幂对应的初等因子分解，指出有限生成和主理想整环假设分别在哪些步骤使用。

\section{有限生成阿贝尔群的分类}
取 $R=\mathbb Z$，由式\eqref{eq:pid-module-structure}得到任意有限生成阿贝尔群为自由阿贝尔群与有限循环群的直和。回访第\ref{ch:abelian}章，将有限情形的初等证明与模论证明对照。通过关系矩阵计算商群，区分无挠部分的秩与有限挠部分的阶。

\section{有理标准形与若尔当标准形}
设 $V$ 为域 $F$ 上的有限维向量空间、$T:V\to V$ 为线性变换。证明对应的 $F[x]$-模为有限生成挠模，由不变因子得到有理标准形（rational canonical form），并读出最小多项式和特征多项式。只有当特征多项式在 $F$ 上分裂时，才在 $F$ 上进一步得到若尔当标准形（Jordan canonical form）；否则需保留有理标准形或扩张基域。

\section{例题、练习与自检}
由一个整数关系矩阵计算有限生成阿贝尔群的不变因子，再由给定最小多项式与特征多项式列出可能的矩阵标准形。训练从各次幂 $\ker(T-\lambda I)^k$ 的维数恢复若尔当块大小，其中 $\lambda$ 为特征值、$I$ 为恒等变换。综合题要求把“求商群”和“判断矩阵相似”分别翻译为模的结构问题。

\backmatter
\chapter{综合复习与后续学习}

\section{贯穿全书的四条线索}
\begin{enumerate}
  \item \textbf{生成与关系。} 从循环群的一个生成元，到商环中规定某个多项式为零，再到模的关系矩阵，比较“给出元素”与“规定元素间的关系”如何共同确定结构。
  \item \textbf{同态与商结构。} 对照群、环、模的核与像，说明各自的商对象需要什么条件，并能实际构造第一同构定理中的映射。
  \item \textbf{作用与对称性。} 从几何对称和置换到域自同构，识别作用对象、固定对象与轨道，理解对称性的限制如何反映原问题的结构。
  \item \textbf{分解与分类。} 对照有限阿贝尔群、唯一分解整环和有限生成模，分别说清楚分解对象、基本因子及唯一性的准确含义。
\end{enumerate}

\section{综合练习的编写规划}
综合题组围绕五项任务展开：分类给定阶数的阿贝尔群并与非阿贝尔例子比较；用理想与中国剩余定理求解联立同余；构造有限域并计算其子域和自同构；由具体多项式求分裂域、伽罗瓦群及中间域；用关系矩阵和不变因子连接商群计算与矩阵标准形。每组均包含至少一个反例或假设辨析题，要求说明结论的适用范围。

\section{后续课程的入口}
交换代数沿理想、局部化与模继续发展，表示论从群在线性空间上的作用出发，代数数论研究数域及其整数环，代数几何连接多项式理想与几何对象。上述方向仅作为后续路线，不将它们的完整理论混入入门章节。自由群、张量积、一般局部化及范畴语言可在完成核心主线后另立专题。

\chapter{参考资料与使用方式}

本书的章节组织以先修依赖与本科自学需要为依据，不对应某一本教材的逐章翻译。以下作者公开讲义用于后续核对群论与域论中的定理条件和证明；其部分内容面向更高层次，宜按主题查阅，不要求初学者通读。

\begin{enumerate}
  \item J. S. Milne, \emph{Group Theory}。作者页面列有群的基本结构、群作用、西罗定理、可解群等内容；用于群论部分的补充核对。\url{https://www.jmilne.org/math/CourseNotes/gt.html}。
  \item J. S. Milne, \emph{Fields and Galois Theory}。作者页面列有域扩张、分裂域、伽罗瓦基本定理、群的计算与应用；用于域论和根式可解性部分的补充核对。\url{https://www.jmilne.org/math/CourseNotes/ft.html}。
\end{enumerate}

\end{document}
